% ============================================================================
% MÓDULO 6. DASHBOARDS INTERACTIVOS CON SHINY
% ============================================================================
\chapter{Dashboards interactivos con Shiny}
\label{cap:shiny}

\begin{objetivos}
\begin{itemize}
  \item Explicar qué hace bueno a un dashboard de BI (audiencia, preguntas,
        KPIs, filtros) y decidir cuándo conviene Shiny y cuándo Power BI o
        Tableau.
  \item Entender la anatomía de una app Shiny (\codigo{ui}, \codigo{server},
        \codigo{shinyApp()}) y ejecutar las tres apps completas que acompañan
        este módulo.
  \item Usar las entradas y salidas más comunes y saber qué función
        \codigo{render...()} corresponde a cada \codigo{...Output()}.
  \item Dominar la \termino{reactividad}: \codigo{reactive()},
        \codigo{observeEvent()}, \codigo{eventReactive()}, \codigo{req()},
        \codigo{validate()} e \codigo{isolate()}, y reconocer sus errores
        típicos.
  \item Organizar la pantalla con \codigo{sidebarLayout()},
        \codigo{fluidRow()}, pestañas y \paquete{shinydashboard}
        (\codigo{valueBox()}, \codigo{box()}, menús).
  \item Agregar tablas interactivas con \paquete{DT}, botones de descarga en
        CSV y Excel, y gráficas interactivas con \paquete{plotly}.
  \item Aplicar buenas prácticas profesionales: separar la lógica de la
        interfaz, módulos, caché, pruebas con \codigo{testServer()} y
        publicación segura.
\end{itemize}
\end{objetivos}

Llevas cinco módulos respondiendo preguntas de negocio con R. Cada vez que la
directora comercial te pide algo (``¿cuánto vendimos en Monterrey en el
segundo trimestre?'', ``¿y solo en Electrónica?'', ``¿y si quitamos el
Outlet?'') tú abres RStudio, cambias un \codigo{filter()}, vuelves a correr
el script y le mandas una gráfica nueva por correo. Funciona, pero no escala:
tú te conviertes en el cuello de botella y ella espera horas por una
respuesta que le tomaría segundos si pudiera mover los filtros por sí misma.

La solución es un \termino{dashboard interactivo}: una página web donde la
directora elige el periodo, la tienda y la categoría, y las cifras y
gráficas se actualizan al instante. Con el paquete \paquete{shiny} puedes
construirlo \emph{con el mismo R que ya sabes}: tus \codigo{filter()},
\codigo{summarise()} y \codigo{ggplot()} siguen siendo el corazón; Shiny solo
añade la ``piel'' (botones, menús, cajas) y el ``sistema nervioso'' (la
reactividad) que conecta lo que el usuario toca con lo que R calcula.

En este módulo construirás tres apps de complejidad creciente, que vienen
completas en la carpeta \archivo{latex/codigo/apps/} del curso: una
calculadora de precios (tu primera app), un explorador de ventas con filtros
y un dashboard ejecutivo con cinco pestañas, indicadores, tablas y botones de
descarga. Al final tendrás la base del dashboard del proyecto integrador
(Módulo~\ref{cap:proyecto}).

% ----------------------------------------------------------------------------
\section{¿Qué es un dashboard de BI y qué lo hace bueno?}
\label{sec:m6-dashboard}

\begin{concepto}{Dashboard de BI}
Un \termino{dashboard} (tablero de control) es una pantalla que reúne los
indicadores más importantes de un negocio para que una audiencia concreta
pueda \emph{monitorear} cómo va y \emph{decidir} qué hacer, de un vistazo.
Un \termino{KPI} (\emph{Key Performance Indicator}, indicador clave de
desempeño) es una cifra que resume algo que importa: ventas, ticket
promedio, clientes activos, margen. Un dashboard \emph{interactivo} además
deja filtrar y profundizar sin pedirle nada al analista.
\end{concepto}

Un dashboard no es ``poner todas las gráficas que tengo en una pantalla''.
Los buenos comparten cinco rasgos:

\begin{enumerate}
  \item \textbf{Tiene una audiencia definida.} El director general quiere
        cinco cifras y una tendencia; el gerente de tienda quiere detalle de
        su sucursal; el analista quiere descargar datos. Un dashboard para
        ``todos'' no le sirve bien a nadie.
  \item \textbf{Responde preguntas concretas.} Antes de programar, escribe
        las 3--5 preguntas que debe contestar (``¿vamos arriba o abajo del
        mes pasado?'', ``¿qué tienda jala hacia abajo?''). Cada elemento de
        la pantalla debe responder a una de ellas.
  \item \textbf{Pone los KPIs arriba.} La vista sigue un patrón en ``Z'': lo
        primero que se lee es la esquina superior izquierda. Ahí van las
        cifras principales en tarjetas grandes; debajo, las tendencias; al
        final, el detalle (tablas).
  \item \textbf{Agrupa los filtros en un solo lugar.} Todos los filtros
        juntos (en una barra lateral o en una fila arriba), con valores por
        defecto sensatos (``todo el año'', ``todas las tiendas'') y que
        afecten a toda la página por igual.
  \item \textbf{Menos es más.} Pocas gráficas, bien elegidas, con un solo
        color salvo que el color signifique algo, ejes con formato de negocio
        (\$, \%, miles) y títulos que digan qué mirar. Lo que aprendiste en el
        Módulo~\ref{cap:visualizacion} aplica aquí multiplicado.
\end{enumerate}

La figura~\ref{fig:m6-anatomia-dashboard} resume esa jerarquía visual. La
usaremos como plantilla para el dashboard ejecutivo de la
sección~\ref{sec:m6-shinydashboard}.

\begin{figure}[H]
\centering
\begin{tikzpicture}[
    font=\sffamily\small,
    zona/.style={draw=gristexto!60, rounded corners=2pt, align=center},
    kpi/.style={zona, fill=primario!12, minimum width=2.55cm, minimum height=1.1cm},
    graf/.style={zona, fill=secundario!10},
    nota/.style={font=\sffamily\footnotesize, text=gristexto, align=left, anchor=west}]
  % marco de la pantalla
  \draw[gristexto, rounded corners=3pt, line width=0.8pt] (0,0) rectangle (11.2,7);
  \fill[primario!85] (0,6.4) rectangle (11.2,7);
  \node[white, anchor=west] at (0.2,6.7) {\textbf{Ventas 2023}};
  % filtros
  \node[zona, fill=acento!12, minimum width=10.8cm, minimum height=0.6cm,
        anchor=north west] at (0.2,6.2)
        {Filtros: \ periodo \quad tienda \quad categoría \quad [Quitar filtros]};
  % kpis
  \foreach \i/\t/\v in {0/Ventas/\$959,993, 1/Transacciones/1{,}000,
                        2/Ticket promedio/\$960, 3/Clientes/198} {
    \node[kpi, anchor=north west] at ({0.2+\i*2.72},5.35)
      {\textbf{\v}\\[-1pt]{\footnotesize \t}};
  }
  % graficas
  \node[graf, minimum width=6.9cm, minimum height=2.2cm, anchor=north west]
       at (0.2,4.0) {Tendencia mensual\\(línea)};
  \node[graf, minimum width=3.7cm, minimum height=2.2cm, anchor=north west]
       at (7.3,4.0) {Composición\\(barras)};
  % tabla
  \node[zona, fill=gray!8, minimum width=10.8cm, minimum height=1.35cm,
        anchor=north west] at (0.2,1.6) {Detalle: tabla con búsqueda y descarga};
  % numeración de lectura
  \foreach \y/\n in {5.9/1, 4.8/2, 2.9/3, 0.9/4} {
    \node[circle, fill=acento, text=white, inner sep=1.5pt, font=\sffamily\scriptsize\bfseries]
      at (11.55,\y) {\n};
  }
  \node[nota] at (11.8,5.9) {Filtros en un\\solo lugar};
  \node[nota] at (11.8,4.8) {KPIs arriba};
  \node[nota] at (11.8,2.9) {Tendencias y\\comparaciones};
  \node[nota] at (11.8,0.9) {Detalle al final};
\end{tikzpicture}
\caption{Jerarquía de un buen dashboard: filtros juntos, KPIs arriba,
tendencias en medio y detalle al final.}
\label{fig:m6-anatomia-dashboard}
\end{figure}

\begin{negocio}{Tres dashboards, tres audiencias}
En una cadena de retail es común tener al menos tres vistas de los mismos
datos: (1) un \emph{dashboard ejecutivo} para la dirección, con 4--6 KPIs,
la tendencia contra el año anterior y el ranking de tiendas; (2) un
\emph{dashboard operativo} para gerentes de tienda, filtrado a su sucursal,
con ventas del día, inventario bajo y desempeño por vendedor; y (3) un
\emph{explorador analítico} para el equipo de BI, con muchos filtros, tablas
detalladas y descarga de datos. Mismo origen, distinto diseño.
\end{negocio}

\subsection{Shiny frente a Power BI y Tableau}

Seguramente conoces herramientas como Power BI o Tableau, que permiten armar
dashboards arrastrando campos con el ratón. ¿Por qué aprender Shiny? No se
trata de que una herramienta sea mejor que otra, sino de elegir la adecuada
para cada caso (tabla~\ref{tab:m6-shiny-vs-bi}).

\begin{table}[H]
\centering
\caption{Cuándo conviene cada herramienta.}
\label{tab:m6-shiny-vs-bi}
\small
\begin{tabularx}{\textwidth}{>{\raggedright\arraybackslash}p{3.1cm}
                             >{\raggedright\arraybackslash}X
                             >{\raggedright\arraybackslash}X}
\toprule
\textbf{Aspecto} & \textbf{Power BI / Tableau} & \textbf{Shiny (R)} \\
\midrule
Curva de aprendizaje & Baja: arrastrar y soltar & Media: hay que programar \\
Tiempo al primer tablero & Minutos & Una o dos horas \\
Lógica compleja (modelos, simulaciones, pronósticos) & Limitada o con
extensiones & Todo lo que R sabe hacer: \paquete{forecast}, modelos de
Machine Learning, optimización \\
Reproducibilidad y control de versiones & Archivos binarios difíciles de
comparar & Código de texto: Git, revisiones, pruebas automáticas \\
Diseño a la medida & Dentro de lo que ofrece la herramienta & Total (HTML,
CSS, JavaScript si hace falta) \\
Costo & Licencia por usuario & Paquete gratuito; el servidor puede ser
gratuito o de pago \\
Ideal para & Reportes corporativos estándar, autoservicio de usuarios de
negocio & Apps analíticas a la medida, prototipos rápidos para quien ya
trabaja en R, herramientas con modelos \\
\bottomrule
\end{tabularx}
\end{table}

\begin{consejo}
Muchas empresas usan ambos: Power BI o Tableau para los reportes estándar que
consume toda la organización, y Shiny para las herramientas analíticas que
necesitan un modelo o una lógica que esas plataformas no cubren (un
simulador de precios, un pronóstico de demanda, un segmentador de clientes).
Saber Shiny te permite construir justo lo que las otras herramientas no
pueden.
\end{consejo}

% ----------------------------------------------------------------------------
\section{Anatomía de una app Shiny}
\label{sec:m6-anatomia}

\subsection{Instalación}

Los paquetes de este módulo se instalan una sola vez:

\begin{rnoejecutar}
install.packages(c("shiny", "shinydashboard", "DT", "plotly", "writexl"))
\end{rnoejecutar}

A lo largo del módulo usaremos estos paquetes. Cárgalos al inicio de tu
script (al cargar \paquete{shinydashboard} y \paquete{DT} verás avisos de que
``enmascaran'' funciones como \codigo{box()} o \codigo{renderDataTable()}; son
normales):

\begin{rcode}
library(shiny)            # apps web con R
library(shinydashboard)   # plantillas de dashboard
library(DT)               # tablas interactivas
library(tidyverse)        # dplyr, ggplot2, readr...
library(lubridate)        # fechas
library(scales)           # formatos de moneda y porcentaje
options(width = 70)
\end{rcode}

\subsection{Las tres piezas: \texorpdfstring{\codigo{ui}, \codigo{server} y \codigo{shinyApp()}}{ui, server y shinyApp()}}

\begin{concepto}{Estructura de toda app Shiny}
Toda app Shiny tiene tres piezas:
\begin{itemize}
  \item \codigo{ui} (\emph{user interface}): describe lo que el usuario
        \textbf{ve}. Es un objeto de R que en realidad es HTML: títulos,
        cajas de entrada, huecos donde irán las gráficas.
  \item \codigo{server}: una función con los argumentos \codigo{input},
        \codigo{output} y \codigo{session} que contiene la \textbf{lógica}:
        lee lo que el usuario eligió (\codigo{input\$...}) y produce
        resultados (\codigo{output\$...}).
  \item \codigo{shinyApp(ui, server)}: une ambas piezas y crea la app.
\end{itemize}
\end{concepto}

Cuando lanzas una app, R arranca un pequeño \termino{servidor web} en tu
computadora. El navegador pide la página, recibe el HTML de la \codigo{ui} y
abre una conexión permanente (un \termino{WebSocket}) con R. Cada vez que
mueves un control, el navegador envía el nuevo valor; R recalcula
\emph{solo} lo que depende de ese valor y devuelve el resultado
(figura~\ref{fig:m6-flujo}).

\begin{figure}[H]
\centering
\begin{tikzpicture}[
    font=\sffamily\small, >={Stealth[length=2.5mm]},
    caja/.style={draw, rounded corners=3pt, line width=0.8pt, align=center,
                 minimum width=4.6cm, minimum height=3.7cm},
    mini/.style={draw=gristexto!70, rounded corners=1.5pt, fill=white,
                 font=\sffamily\footnotesize, minimum width=3.6cm, align=center}]
  \node[caja, draw=secundario, fill=secundario!6] (nav) at (0,0) {};
  \node[font=\sffamily\bfseries, text=secundario, anchor=north]
       at (nav.north) {Navegador (lo que ve el usuario)};
  \node[mini] (in1) at (0,0.55) {Precio: [ 1200 ]};
  \node[mini] (in2) at (0,-0.15) {Descuento: \textbf{-----o-----} 25\%};
  \node[mini, fill=exito!8] (out1) at (0,-1.05) {Precio final: \$900};
  \node[caja, draw=primario, fill=primario!5] (r) at (8.2,0) {};
  \node[font=\sffamily\bfseries, text=primario, anchor=north]
       at (r.north) {Sesión de R (server)};
  \node[mini, font=\ttfamily\footnotesize] (inr) at (8.2,0.4)
       {input\$precio\\input\$descuento};
  \node[mini, font=\ttfamily\footnotesize, fill=primario!8] (calc) at (8.2,-1.05)
       {output\$precio\_final <-\\renderText(\{...\})};
  \draw[->, thick, acento] (in1.east) to[out=10, in=170]
       node[above, font=\sffamily\scriptsize, text=acento] {1. el usuario cambia un valor} (inr.west);
  \draw[->, thick, primario] (inr.south) -- node[right, font=\sffamily\scriptsize] {2. R recalcula} (calc.north);
  \draw[->, thick, exito] (calc.west) --
       node[below, font=\sffamily\scriptsize, text=exito, align=center] {3. envía el resultado\\(texto, imagen, HTML)} (out1.east);
  \node[font=\sffamily\scriptsize, text=gristexto] at (4.1,-2.35)
       {Conexión permanente (WebSocket) mientras la pestaña esté abierta};
\end{tikzpicture}
\caption{Flujo de una app Shiny: el navegador envía entradas, R recalcula y
devuelve salidas.}
\label{fig:m6-flujo}
\end{figure}

\subsection{Tu primera app paso a paso}
\label{sec:m6-primera-app}

\begin{ejemplo}{Calculadora de precio con descuento}
El área comercial quiere que los vendedores calculen el precio final de un
producto con descuento sin abrir Excel. La app debe tener dos entradas
(precio de lista y porcentaje de descuento) y mostrar el precio final. Está
en \archivo{latex/codigo/apps/01-primera-app/app.R}.
\end{ejemplo}

Este es el archivo completo. Léelo con calma: son apenas 30 líneas de código
útil y contienen todo lo esencial de Shiny.

\begin{rnoejecutar}[title=latex/codigo/apps/01-primera-app/app.R]
# 1. Paquete -------------------------------------------------------------
library(shiny)          # el motor de las apps web con R

# 2. Interfaz (ui): lo que el usuario VE en el navegador -----------------
ui <- fluidPage(                       # página que se adapta al ancho
  titlePanel("Calculadora de precio con descuento"),   # título

  # Entrada 1: una caja para escribir un número (el precio de lista)
  numericInput(inputId = "precio", label = "Precio de lista ($):",
               value = 500, min = 0, step = 10),

  # Entrada 2: un deslizador para elegir el descuento
  sliderInput(inputId = "descuento", label = "Descuento (%):",
              min = 0, max = 50, value = 10, step = 5),

  # Salida: un hueco donde el servidor escribirá el resultado
  h3(textOutput(outputId = "precio_final"))
)

# 3. Servidor (server): la lógica que CORRE en R -------------------------
server <- function(input, output, session) {
  # renderText() se vuelve a ejecutar SOLO cuando cambian las entradas
  # que usa (input$precio o input$descuento): eso es la reactividad
  output$precio_final <- renderText({
    final <- input$precio * (1 - input$descuento / 100)
    paste0("Precio final: ", scales::dollar(final))
  })
}

# 4. Unir interfaz y servidor: esto crea (y lanza) la app ----------------
shinyApp(ui = ui, server = server)
\end{rnoejecutar}

Vamos por partes:

\begin{itemize}
  \item \textbf{Línea 2:} \codigo{library(shiny)} carga el paquete. Toda app
        empieza así.
  \item \textbf{Líneas 5--18, la \codigo{ui}:} \codigo{fluidPage()} crea una
        página que se adapta al ancho de la pantalla (de computadora o de
        celular). Dentro van, separados \emph{por comas}, sus elementos: un
        título, dos \termino{entradas} (\emph{inputs}) y una
        \termino{salida} (\emph{output}).
  \item \textbf{El \codigo{inputId}:} cada entrada tiene un identificador
        único (\codigo{"precio"}, \codigo{"descuento"}). Es el nombre con el
        que el servidor lee su valor: \codigo{input\$precio}.
  \item \textbf{\codigo{textOutput("precio\_final")}:} reserva un hueco en
        la página. No calcula nada; solo dice ``aquí irá un texto llamado
        \codigo{precio\_final}''. Lo envolvemos en \codigo{h3()} para que se
        vea grande (\codigo{h3} es un encabezado de HTML).
  \item \textbf{Líneas 21--28, el \codigo{server}:} es una función. Dentro
        asignamos \codigo{output\$precio\_final} con \codigo{renderText()}.
        El código entre llaves lee \codigo{input\$precio} e
        \codigo{input\$descuento}, calcula y devuelve un texto.
  \item \textbf{La magia:} no escribimos ningún ``cuando el usuario mueva el
        deslizador, recalcula''. Shiny detecta solo que
        \codigo{renderText()} usó esas dos entradas y lo vuelve a ejecutar
        cada vez que alguna cambia. Eso es la \termino{reactividad}.
  \item \textbf{Línea 31:} \codigo{shinyApp()} une las piezas. Si ejecutas
        el archivo, la app se abre.
\end{itemize}

La \codigo{ui} es solo un objeto de R que genera HTML. Puedes construir una
entrada en la consola e imprimirla para ver lo que el navegador recibirá.
Es un gran recurso para entender (y depurar) Shiny:

\begin{rcode}
# Una entrada numérica, igual a la de la app, impresa en la consola
numericInput(inputId = "precio", label = "Precio de lista ($):",
             value = 500, min = 0, step = 10)
\end{rcode}
\begin{rsalida}
<div class="form-group shiny-input-container">
  <label class="control-label" id="precio-label" for="precio">Precio de lista ($):</label>
  <input id="precio" type="number" class="shiny-input-number form-control" value="500" min="0" step="10"/>
</div>
\end{rsalida}

Observa el \codigo{id="precio"}: el \codigo{inputId} se convierte en el
identificador del elemento HTML, y por eso \textbf{debe ser único} en toda la
app. Las salidas son todavía más simples, un contenedor vacío con una clase
especial que Shiny sabe llenar:

\begin{rcode}
textOutput(outputId = "precio_final")
plotOutput(outputId = "grafica", height = "300px")
\end{rcode}
\begin{rsalida}
<div id="precio_final" class="shiny-text-output"></div>
<div class="shiny-plot-output html-fill-item" id="grafica" style="width:100%;height:300px;"></div>
\end{rsalida}

\begin{figure}[H]
\centering
\begin{tikzpicture}[font=\sffamily\small]
  \draw[gristexto, rounded corners=3pt, line width=0.8pt] (0,0) rectangle (9,4.6);
  \fill[gray!12] (0,4.1) rectangle (9,4.6);
  \foreach \x in {0.25,0.5,0.75} \fill[gray!45] (\x,4.35) circle (0.07);
  \node[anchor=west, font=\sffamily\footnotesize, text=gristexto] at (1,4.35) {http://127.0.0.1:XXXX};
  \node[anchor=west, font=\sffamily\Large] at (0.3,3.55) {Calculadora de precio con descuento};
  \node[anchor=west] at (0.3,2.85) {Precio de lista (\$):};
  \draw[gristexto!80, fill=white] (0.3,2.1) rectangle (4.3,2.6);
  \node[anchor=west] at (0.4,2.35) {500};
  \node[anchor=west] at (0.3,1.7) {Descuento (\%):};
  \draw[line width=2pt, secundario!50] (0.4,1.2) -- (4.2,1.2);
  \draw[line width=2pt, secundario] (0.4,1.2) -- (1.16,1.2);
  \fill[white, draw=secundario] (1.16,1.2) circle (0.12);
  \node[font=\sffamily\scriptsize, text=gristexto] at (0.4,0.9) {0};
  \node[font=\sffamily\scriptsize, text=gristexto] at (4.2,0.9) {50};
  \node[font=\sffamily\scriptsize, text=secundario] at (1.16,1.5) {10};
  \node[anchor=west, font=\sffamily\large\bfseries] at (0.3,0.4) {Precio final: \$450};
  \draw[<-, acento, thick] (4.4,2.35) -- (5.6,2.35) node[right, font=\sffamily\footnotesize] {\codigo{numericInput("precio")}};
  \draw[<-, acento, thick] (4.4,1.2) -- (5.6,1.2) node[right, font=\sffamily\footnotesize] {\codigo{sliderInput("descuento")}};
  \draw[<-, exito, thick] (3.6,0.4) -- (5.6,0.4) node[right, font=\sffamily\footnotesize] {\codigo{textOutput("precio\_final")}};
\end{tikzpicture}
\caption{Maqueta de la app 01 tal como se ve en el navegador.}
\label{fig:m6-maqueta-app01}
\end{figure}

\subsection{Cómo ejecutar (y detener) una app}
\label{sec:m6-ejecutar}

Tienes tres formas de lanzar una app guardada en un archivo \archivo{app.R}:

\begin{enumerate}
  \item Abrir \archivo{app.R} en RStudio y pulsar el botón
        \menu{Run App} (arriba a la derecha del editor). Con la flechita del
        botón eliges si se abre en el visor de RStudio o en tu navegador.
  \item Desde la consola, con el directorio de trabajo en la carpeta del
        curso, pasar a \codigo{runApp()} la \emph{carpeta} que contiene el
        \archivo{app.R}:
\end{enumerate}

\begin{rnoejecutar}
# Directorio de trabajo: la carpeta del curso (la que contiene datasets/)
getwd()
shiny::runApp("latex/codigo/apps/01-primera-app")
shiny::runApp("latex/codigo/apps/02-explorador-ventas")
shiny::runApp("latex/codigo/apps/03-dashboard-ejecutivo")
\end{rnoejecutar}

\begin{enumerate}[start=3]
  \item Si defines \codigo{ui} y \codigo{server} en tu script, ejecutar
        \codigo{shinyApp(ui, server)} en la consola también la abre.
\end{enumerate}

Mientras la app está abierta, la consola de R queda \emph{ocupada} (verás
\codigo{Listening on http://127.0.0.1:XXXX}) y no acepta otros comandos. Para
detenerla pulsa \tecla{Esc} en la consola, el botón rojo \menu{Stop} del
panel de consola, o cierra la ventana de la app.

\begin{advertencia}
Todo lo que escribe \codigo{print()} o \codigo{message()} dentro del
\codigo{server} aparece en la \emph{consola de R}, no en la página. Úsalo
para depurar: si no sabes qué valor tiene una entrada, escribe
\codigo{print(input\$descuento)} dentro del \codigo{renderText()} y mira la
consola mientras mueves el control.
\end{advertencia}

\subsection{El truco para que las apps encuentren \texorpdfstring{\archivo{datasets/}}{datasets/}}
\label{sec:m6-rutas}

Aquí hay una trampa en la que cae casi todo el mundo. Cuando ejecutas
\codigo{runApp("latex/codigo/apps/02-explorador-ventas")}, Shiny
\textbf{cambia el directorio de trabajo a la carpeta de la app} mientras
esta corre. Por eso, dentro de la app, la ruta relativa
\archivo{datasets/transacciones.csv} ya no existe y \codigo{read\_csv()}
falla con un mensaje como este (recortado):

\begin{textocodigo}
Error: 'datasets/transacciones.csv' does not exist in current working
directory ('C:/Users/tu_usuario/curso-bi-rstudio/latex/codigo/apps/02-explorador-ventas').
\end{textocodigo}

Hay varias soluciones: copiar los datos dentro de la carpeta de la app (lo
normal al publicarla), usar una ruta absoluta (frágil: solo funciona en tu
computadora) o, la que usan las apps de este libro, \textbf{buscar la carpeta
\archivo{datasets/} subiendo directorios}: primero en la carpeta actual,
luego en la de arriba (\archivo{..}), luego en la de más arriba, hasta
encontrarla. Así la app funciona tanto si la lanzas desde la carpeta del
curso como si pulsas \menu{Run App} con el archivo abierto.

\begin{rcode}
# Busca datasets/transacciones.csv subiendo carpetas desde 'desde'
buscar_datasets <- function(desde = ".", max_niveles = 6) {
  carpeta <- desde
  for (i in 0:max_niveles) {
    candidato <- file.path(carpeta, "datasets", "transacciones.csv")
    if (file.exists(candidato)) {
      return(file.path(carpeta, "datasets"))   # ¡encontrada!
    }
    carpeta <- file.path(carpeta, "..")        # sube un nivel
  }
  stop("No encontré datasets/transacciones.csv. Ejecuta primero ",
       "source('datasets/generar_datasets.R') en la carpeta del curso.",
       call. = FALSE)
}
\end{rcode}

Pruébala. Desde la carpeta del curso la encuentra a la primera; desde la
carpeta de la app 02 tiene que subir cuatro niveles:

\begin{rcode}
buscar_datasets()                    # desde la carpeta del curso
carpeta_app <- "latex/codigo/apps/02-explorador-ventas"
buscar_datasets(desde = carpeta_app) # como la vería la app
file.exists(file.path(buscar_datasets(carpeta_app), "tiendas.csv"))
\end{rcode}
\begin{rsalida}
[1] "./datasets"
[1] "latex/codigo/apps/02-explorador-ventas/../../../../datasets"
[1] TRUE
\end{rsalida}

La segunda ruta se lee ``desde la carpeta de la app, sube cuatro veces y
entra a \archivo{datasets}''. Si la función no encuentra los datos, se
detiene con un mensaje en español que dice exactamente qué hacer, mucho más
útil que el error genérico de \codigo{read\_csv()}.

\begin{hazlotu}
Ejecuta \codigo{shiny::runApp("latex/codigo/apps/01-primera-app")}, cambia
el precio a 1,200 y el descuento a 25\%. ¿Qué precio final ves? (Debe ser
\$900). Después detén la app con \tecla{Esc}.
\end{hazlotu}

% ----------------------------------------------------------------------------
\section{Entradas y salidas}
\label{sec:m6-entradas-salidas}

\subsection{Las entradas más útiles en BI}

Una \termino{entrada} (\emph{input}) es cualquier control con el que el
usuario le dice algo a la app. Todas las funciones de entrada tienen la misma
forma: \codigo{xxxInput(inputId, label, ...)}. El primer argumento es el
identificador (lo que usarás como \codigo{input\$inputId}), el segundo la
etiqueta visible y el resto depende del control. La
tabla~\ref{tab:m6-entradas} resume las que más usarás y, muy importante,
\emph{qué tipo de valor llega} al servidor.

\begin{table}[H]
\centering
\caption{Entradas frecuentes en dashboards y el valor que recibe el servidor.}
\label{tab:m6-entradas}
\small
\begin{tabularx}{\textwidth}{>{\raggedright\arraybackslash}p{3.6cm}
                             >{\raggedright\arraybackslash}X
                             >{\raggedright\arraybackslash}p{4.3cm}}
\toprule
\textbf{Función} & \textbf{Uso típico} & \textbf{Valor en \codigo{input\$id}} \\
\midrule
\codigo{selectInput()} & Elegir una tienda, un año, una métrica & Texto
(\codigo{"Sucursal Centro"}) \\
\codigo{selectizeInput(..., multiple = TRUE)} & Elegir varias categorías,
con buscador & Vector de texto, o \codigo{NULL} si no hay nada elegido \\
\codigo{dateRangeInput()} & Periodo de análisis & Vector de 2 fechas
(\codigo{Date}) \\
\codigo{dateInput()} & Una fecha (corte, día) & Una fecha \\
\codigo{sliderInput()} & Top N, rango de precios, un porcentaje & Número (o
2 números si \codigo{value} tiene 2) \\
\codigo{numericInput()} & Metas, presupuestos, precios & Número (\codigo{NA}
si el usuario lo borra) \\
\codigo{checkboxGroupInput()} & Varias opciones visibles a la vez (métodos
de pago) & Vector de texto, o \codigo{NULL} \\
\codigo{radioButtons()} & Una opción entre pocas (ventas / unidades) &
Texto \\
\codigo{checkboxInput()} & Sí/no (``incluir outlet'') & \codigo{TRUE} o
\codigo{FALSE} \\
\codigo{textInput()} & Buscar un cliente por nombre & Texto \\
\codigo{actionButton()} & ``Calcular'', ``Actualizar'', ``Quitar filtros'' &
Contador: 0, 1, 2\ldots{} (veces que se pulsó) \\
\bottomrule
\end{tabularx}
\end{table}

\begin{advertencia}
Presta atención a las entradas \textbf{múltiples}: cuando el usuario no ha
elegido nada, el valor es \codigo{NULL}, no un vector vacío de texto ni
\codigo{"Todas"}. Si filtras con \codigo{categoria \%in\% input\$categorias}
sin revisar eso, te quedarás sin filas. La convención de este libro es:
\emph{vacío significa ``todas''}, y lo comprobamos con
\codigo{length(input\$categorias) > 0}.
\end{advertencia}

Veamos cómo se construyen, con datos reales del curso. Las opciones de los
menús casi nunca se escriben a mano: se sacan de los datos, así la app no se
desactualiza cuando abre una tienda nueva.

\begin{rcode}
# Datos para los ejemplos de esta sección
transacciones <- read_csv("datasets/transacciones.csv",
                          show_col_types = FALSE)
tiendas <- read_csv("datasets/tiendas.csv", show_col_types = FALSE)

# Las opciones salen de los datos, ordenadas
opciones_tienda <- sort(tiendas$nombre_tienda)
opciones_pago   <- sort(unique(transacciones$metodo_pago))
opciones_pago
\end{rcode}
\begin{rsalida}
[1] "Efectivo"        "Tarjeta Crédito" "Tarjeta Débito" 
[4] "Transferencia"  
\end{rsalida}

Con esos vectores armamos un ``catálogo'' con las entradas de la
tabla~\ref{tab:m6-entradas}. Fíjate en las comas entre elementos y en los
argumentos que controlan el valor inicial:

\begin{rcode}
ui_catalogo <- fluidPage(
  selectInput("tienda", "Tienda:", choices = opciones_tienda,
              selected = "Sucursal Centro"),
  selectizeInput("pagos", "Métodos de pago:", choices = opciones_pago,
                 multiple = TRUE, options = list(placeholder = "Todos")),
  dateRangeInput("fechas", "Periodo:", start = "2023-01-01",
                 end = "2023-12-31", language = "es", separator = " a ",
                 format = "dd/mm/yyyy"),
  sliderInput("top", "Top N:", min = 5, max = 20, value = 10, step = 5),
  numericInput("meta", "Meta mensual ($):", value = 80000, step = 5000),
  checkboxGroupInput("trimestres", "Trimestres:",
                     choices = c("Q1", "Q2", "Q3", "Q4"),
                     selected = c("Q1", "Q2", "Q3", "Q4"), inline = TRUE),
  radioButtons("metrica", "Medir por:",
               choices = c("Ventas ($)" = "ventas",
                           "Unidades" = "unidades")),
  actionButton("calcular", "Calcular", icon = icon("play"))
)
\end{rcode}

En \codigo{radioButtons()} usamos un vector \emph{con nombres}:
\codigo{"Ventas (\$)" = "ventas"}. El nombre es lo que ve el usuario y el
valor (\codigo{"ventas"}) es lo que llega a \codigo{input\$metrica}. Así
puedes cambiar la etiqueta sin tocar la lógica. Imprime una de las entradas
para ver cómo quedan las opciones en HTML:

\begin{rcode}
radioButtons("metrica", "Medir por:",
             choices = c("Ventas ($)" = "ventas", "Unidades" = "unidades"))
\end{rcode}
\begin{rsalida}
<div id="metrica" class="form-group shiny-input-radiogroup shiny-input-container" role="radiogroup" aria-labelledby="metrica-label">
  <label class="control-label" id="metrica-label" for="metrica">Medir por:</label>
  <div class="shiny-options-group">
    <div class="radio">
      <label>
        <input type="radio" name="metrica" value="ventas" checked="checked"/>
        <span>Ventas ($)</span>
      </label>
    </div>
    <div class="radio">
      <label>
        <input type="radio" name="metrica" value="unidades"/>
        <span>Unidades</span>
      </label>
    </div>
  </div>
</div>
\end{rsalida}

El atributo \codigo{value="ventas"} es lo que viaja al servidor; el texto
entre \codigo{<span>} es solo la etiqueta.

\begin{hazlotu}
En la consola, imprime \codigo{selectInput("tienda", "Tienda:",
choices = opciones\_tienda)} y localiza cuál opción aparece con la palabra
\codigo{selected}. ¿Qué pasa si agregas \codigo{selected = "Sucursal Sur"}?
\end{hazlotu}

\subsection{Salidas y sus funciones \texorpdfstring{\codigo{render}}{render}}

Una \termino{salida} siempre tiene dos mitades que deben coincidir: en la
\codigo{ui}, una función \codigo{xxxOutput("id")} que reserva el lugar; en el
\codigo{server}, una asignación \codigo{output\$id <- renderXxx(\{...\})} que
produce el contenido. Si usas un \codigo{plotOutput()} con un
\codigo{renderTable()}, o escribes el id distinto en cada lado, la salida
simplemente no aparece (sin ningún error). La
tabla~\ref{tab:m6-salidas} es la que querrás tener a mano.

\begin{table}[H]
\centering
\caption{Parejas de funciones de salida (\codigo{ui}) y de render (\codigo{server}).}
\label{tab:m6-salidas}
\small
\begin{tabularx}{\textwidth}{>{\raggedright\arraybackslash}p{3.9cm}
                             >{\raggedright\arraybackslash}p{3.7cm}
                             >{\raggedright\arraybackslash}X}
\toprule
\textbf{En la \codigo{ui}} & \textbf{En el \codigo{server}} &
\textbf{Qué debe devolver el código} \\
\midrule
\codigo{textOutput()} & \codigo{renderText()} & Un texto (un KPI, un
mensaje) \\
\codigo{verbatimTextOutput()} & \codigo{renderPrint()} & Lo que imprimiría
la consola (\codigo{summary()}, \codigo{str()}) \\
\codigo{plotOutput()} & \codigo{renderPlot()} & Una gráfica (un objeto
\codigo{ggplot}) \\
\codigo{tableOutput()} & \codigo{renderTable()} & Un data frame pequeño
(tabla estática) \\
\codigo{DT::DTOutput()} & \codigo{DT::renderDT()} & Un data frame o un
\codigo{datatable()} (tabla con búsqueda, orden y páginas) \\
\codigo{valueBoxOutput()} & \codigo{renderValueBox()} & Un
\codigo{valueBox()} de \paquete{shinydashboard} \\
\codigo{infoBoxOutput()} & \codigo{renderInfoBox()} & Un \codigo{infoBox()} \\
\codigo{plotly::plotlyOutput()} & \codigo{plotly::renderPlotly()} & Una
gráfica \paquete{plotly} (por ejemplo, \codigo{ggplotly(p)}) \\
\codigo{uiOutput()} & \codigo{renderUI()} & Elementos de interfaz creados
al vuelo \\
\codigo{downloadButton()} & \codigo{downloadHandler()} & Un archivo para
descargar (sección~\ref{sec:m6-descargas}) \\
\bottomrule
\end{tabularx}
\end{table}

\subsection{Probar el servidor sin navegador: \texorpdfstring{\codigo{testServer()}}{testServer()}}

Como no podemos ``hacer clic'' desde un script, usaremos a lo largo del
módulo una herramienta profesional: \codigo{testServer()}. Recibe la función
\codigo{server} (o la carpeta de una app) y un bloque de código donde puedes
simular al usuario con \codigo{session\$setInputs()} y leer las salidas con
\codigo{output\$id}, todo desde la consola. Es la forma de comprobar que la
lógica de tu app funciona sin abrir nada.

\begin{ejemplo}{Resumen de ventas por tienda y método de pago}
Queremos una mini app donde el usuario elija una tienda y, opcionalmente,
uno o varios métodos de pago, y vea cuántas transacciones y cuánto se
vendió. Aprovecharemos las entradas del catálogo anterior.
\end{ejemplo}

\begin{rcode}
server_catalogo <- function(input, output, session) {
  output$resumen <- renderText({
    datos <- transacciones %>%
      left_join(tiendas, by = "tienda_id") %>%
      filter(nombre_tienda == input$tienda)
    if (length(input$pagos) > 0) {          # vacío = todos los métodos
      datos <- filter(datos, metodo_pago %in% input$pagos)
    }
    paste0(input$tienda, ": ", nrow(datos), " transacciones, ",
           dollar(sum(datos$total_transaccion), accuracy = 1))
  })
}

testServer(server_catalogo, {
  session$setInputs(tienda = "Sucursal Centro", pagos = NULL)
  print(output$resumen)                      # todos los métodos
  session$setInputs(pagos = c("Efectivo", "Transferencia"))
  print(output$resumen)                      # solo dos métodos
})
\end{rcode}
\begin{rsalida}
[1] "Sucursal Centro: 87 transacciones, $88,414"
[1] "Sucursal Centro: 36 transacciones, $30,727"
\end{rsalida}

La Sucursal Centro tuvo 87 transacciones en el año, por \$88,414; si nos
quedamos solo con efectivo y transferencia, bajan a 36. Fíjate en que la segunda llamada a
\codigo{setInputs()} solo cambia \codigo{pagos}: la tienda conserva su valor,
igual que en el navegador. Para ver la app en acción bastaría con:

\begin{rnoejecutar}
shinyApp(ui = fluidPage(ui_catalogo, h3(textOutput("resumen"))),
         server = server_catalogo)
\end{rnoejecutar}

% ----------------------------------------------------------------------------
\section{Reactividad: el concepto más importante}
\label{sec:m6-reactividad}

\begin{concepto}{Reactividad}
La \termino{reactividad} es el mecanismo por el que Shiny sabe \emph{qué}
recalcular y \emph{cuándo}. Funciona como una hoja de cálculo: si la celda
C1 contiene \codigo{=A1*B1}, al cambiar A1 se recalcula C1 sola. En Shiny
hay tres tipos de piezas:
\begin{itemize}
  \item \textbf{Fuentes reactivas}: las entradas (\codigo{input\$...}). Cambian
        cuando el usuario actúa.
  \item \textbf{Conductores}: expresiones intermedias creadas con
        \codigo{reactive()}. Calculan algo a partir de las fuentes y guardan
        el resultado (\emph{caché}) hasta que una fuente cambie.
  \item \textbf{Consumidores (puntos finales)}: las salidas
        (\codigo{output\$...}) y los observadores (\codigo{observe()},
        \codigo{observeEvent()}). Son los que realmente ``hacen'' algo.
\end{itemize}
\end{concepto}

Shiny construye automáticamente un \termino{grafo reactivo}: un mapa de qué
depende de qué. La figura~\ref{fig:m6-grafo} muestra el de la app 02 que
construiremos en la sección~\ref{sec:m6-app02}. Cuando el usuario cambia la
tienda, Shiny marca como ``desactualizado'' todo lo que está río abajo de
\codigo{input\$tienda} (el filtrado y las tres salidas) y lo recalcula. Si
en cambio mueve el deslizador \codigo{n\_top}, solo se recalcula la tabla:
ni el filtrado ni la gráfica se tocan.

\begin{figure}[H]
\centering
\begin{tikzpicture}[
    font=\sffamily\small, >={Stealth[length=2.2mm]},
    fuente/.style={draw=acento, fill=acento!10, rounded corners=6pt,
                   minimum width=2.9cm, minimum height=0.65cm, font=\ttfamily\footnotesize},
    conductor/.style={draw=secundario, fill=secundario!12, rounded corners=2pt,
                      minimum width=3.2cm, minimum height=0.9cm, font=\ttfamily\footnotesize},
    salida/.style={draw=primario, fill=primario!10, minimum width=3.1cm,
                   minimum height=0.65cm, font=\ttfamily\footnotesize},
    flecha/.style={->, gristexto, thick}]
  \node[fuente] (tienda) at (0,3) {input\$tienda};
  \node[fuente] (cat) at (0,2) {input\$categorias};
  \node[fuente] (fechas) at (0,1) {input\$fechas};
  \node[fuente] (metrica) at (0,-0.1) {input\$metrica};
  \node[fuente] (top) at (0,-1.1) {input\$n\_top};
  \node[conductor] (datos) at (5,2) {datos\_filtrados()};
  \node[salida] (resumen) at (10,3) {output\$resumen};
  \node[salida] (grafica) at (10,1.3) {output\$grafica};
  \node[salida] (tabla) at (10,-0.4) {output\$tabla\_top};
  \foreach \n in {tienda, cat, fechas} \draw[flecha] (\n.east) -- (datos.west);
  \draw[flecha] (datos.east) -- (resumen.west);
  \draw[flecha] (datos.east) -- (grafica.west);
  \draw[flecha] (datos.east) -- (tabla.west);
  \draw[flecha] (metrica.east) to[out=0, in=200] (grafica.west);
  \draw[flecha] (top.east) -- (tabla.west);
  \node[font=\sffamily\bfseries\footnotesize, text=acento] at (0,3.8) {Fuentes (entradas)};
  \node[font=\sffamily\bfseries\footnotesize, text=secundario] at (5,3.8) {Conductor (reactive)};
  \node[font=\sffamily\bfseries\footnotesize, text=primario] at (10,3.8) {Consumidores (salidas)};
\end{tikzpicture}
\caption{Grafo reactivo de la app 02: las entradas alimentan una expresión
reactiva que comparten las tres salidas.}
\label{fig:m6-grafo}
\end{figure}

\subsection{\texorpdfstring{\codigo{reactive()}}{reactive()}: calcula una vez, úsalo muchas}

En un dashboard, casi todas las salidas parten de los \emph{mismos datos
filtrados}. Si cada \codigo{render...()} filtrara por su cuenta, repetirías
código y harías trabajar a R varias veces. La solución es una
\termino{expresión reactiva}: \codigo{reactive(\{...\})} guarda el
resultado y solo lo recalcula si cambia alguna entrada que usó.

\begin{ejemplo}{Dos KPIs a partir de un solo filtrado}
Queremos mostrar el número de transacciones y el total vendido de la tienda
elegida. Para \emph{demostrar} que el filtrado ocurre una sola vez por
tienda, el reactivo deja un mensaje en la consola cada vez que se ejecuta.
\end{ejemplo}

\begin{rcode}
server_kpis <- function(input, output, session) {
  ventas_tienda <- reactive({             # el conductor
    message("  (filtrando la tienda ", input$tienda, ")")  # rastro
    filter(transacciones, tienda_id == input$tienda)
  })

  output$n_ventas <- renderText(nrow(ventas_tienda()))
  output$total    <- renderText(
    dollar(sum(ventas_tienda()$total_transaccion), accuracy = 1))
}

testServer(server_kpis, {
  session$setInputs(tienda = 3)
  cat("Tienda 3:", output$n_ventas, "ventas,", output$total, "\n")
  session$setInputs(tienda = 8)
  cat("Tienda 8:", output$n_ventas, "ventas,", output$total, "\n")
})
\end{rcode}
\begin{rsalida}
  (filtrando la tienda 3)
Tienda 3: 109 ventas, $105,305 
  (filtrando la tienda 8)
Tienda 8: 106 ventas, $114,377 
\end{rsalida}

Leímos dos salidas por cada tienda, pero el mensaje ``filtrando'' aparece
una sola vez por tienda: la segunda salida reutilizó el resultado guardado.
La tienda 3 (Sucursal Sur) tuvo 109 ventas por \$105,305 y la 8 (Outlet),
106 ventas por \$114,377. Con 1,000 filas no
se nota; con millones de registros o una consulta a una base de datos
(Módulo~\ref{cap:bases-datos}) es la diferencia entre una app ágil y una
desesperante.

\begin{advertencia}
Un reactivo \textbf{se llama como función, con paréntesis}:
\codigo{ventas\_tienda()}. Si escribes \codigo{ventas\_tienda\$total} (sin
paréntesis) obtienes uno de los errores más famosos de Shiny:
\codigo{object of type 'closure' is not subsettable} (en español,
\codigo{objeto de tipo 'closure' no es subconjunto}). ``Closure'' es como R
llama a las funciones: le estás pidiendo una columna a una función. Lo
verás en acción en la sección~\ref{sec:m6-errores}.
\end{advertencia}

\subsection{\texorpdfstring{\codigo{observe()} y \codigo{observeEvent()}}{observe() y observeEvent()}: efectos secundarios}

Las salidas \emph{producen algo que se muestra}. A veces necesitas
\emph{hacer} algo sin mostrar nada: escribir en un registro, mostrar una
notificación, guardar un archivo o actualizar otra entrada. Para eso están
los \termino{observadores}:

\begin{itemize}
  \item \codigo{observe(\{...\})} se ejecuta cada vez que cambia
        \emph{cualquier} entrada que use dentro.
  \item \codigo{observeEvent(evento, \{...\})} se ejecuta solo cuando cambia
        \codigo{evento} (típicamente un botón). Es el más usado porque deja
        claro qué lo dispara.
\end{itemize}

\begin{ejemplo}{Botón ``Quitar filtros'' y registro de uso}
El dashboard ejecutivo tendrá un botón que regresa los filtros a sus
valores iniciales, y al equipo de BI le interesa saber cuántas veces se
usa. Contaremos los clics con un \codigo{reactiveVal()}, un valor reactivo
que tú mismo puedes leer (\codigo{clics()}) y modificar
(\codigo{clics(nuevo\_valor)}).
\end{ejemplo}

\begin{rcode}
server_limpiar <- function(input, output, session) {
  clics <- reactiveVal(0)                     # valor reactivo propio

  observeEvent(input$limpiar, {               # solo al pulsar el botón
    clics(clics() + 1)                        # suma uno
    updateSelectInput(session, "tienda", selected = "Todas")
    message("Filtros reiniciados (clic número ", clics(), ")")
  })

  output$uso <- renderText(paste("Clics en 'Quitar filtros':", clics()))
}

testServer(server_limpiar, {
  session$setInputs(limpiar = 1)              # primer clic
  session$setInputs(limpiar = 2)              # segundo clic
  print(output$uso)
})
\end{rcode}
\begin{rsalida}
Filtros reiniciados (clic número 1)
Filtros reiniciados (clic número 2)
[1] "Clics en 'Quitar filtros': 2"
\end{rsalida}

\codigo{updateSelectInput()} (y sus hermanas \codigo{updateDateRangeInput()},
\codigo{updateSelectizeInput()}\ldots) cambia desde el servidor el valor de
una entrada que ya está en la página. El mensaje aparece en la consola de R,
no en la app: es el lugar ideal para dejar un rastro de lo que ocurre.

\subsection{\texorpdfstring{\codigo{eventReactive()}}{eventReactive()}: calcular solo al pulsar un botón}

Por defecto, Shiny recalcula en cuanto cambias cualquier filtro. Si un
cálculo es pesado (un modelo, un pronóstico, una consulta grande) es mejor
que el usuario ajuste todos los filtros con calma y luego pulse
``Calcular''. \codigo{eventReactive(evento, \{...\})} es un
\codigo{reactive()} que \emph{solo} se actualiza cuando ocurre el evento.

\begin{rcode}
server_meta <- function(input, output, session) {
  # Se recalcula SOLO cuando se pulsa input$calcular
  proyeccion <- eventReactive(input$calcular, {
    ventas_mes <- sum(transacciones$total_transaccion) / 12
    ventas_mes * (1 + input$crecimiento / 100)
  })
  output$proyeccion <- renderText(
    paste("Venta mensual proyectada:", dollar(proyeccion(), accuracy = 1)))
}

testServer(server_meta, {
  session$setInputs(crecimiento = 10)     # el usuario mueve el control
  resultado <- tryCatch(output$proyeccion,
                        error = function(e) "(vacío: aún no se pulsa)")
  print(resultado)
  session$setInputs(calcular = 1)         # ahora sí pulsa el botón
  print(output$proyeccion)
  session$setInputs(crecimiento = 20)     # cambia, pero no pulsa
  print(output$proyeccion)                # sigue el valor anterior
})
\end{rcode}
\begin{rsalida}
[1] "(vacío: aún no se pulsa)"
[1] "Venta mensual proyectada: $87,999"
[1] "Venta mensual proyectada: $87,999"
\end{rsalida}

Antes del primer clic la salida queda vacía; al pulsar se calcula con 10\%
(la venta mensual promedio de \$79,999 más 10\% da \$87,999); y cambiar a
20\% \emph{sin} pulsar no altera nada. Es exactamente lo que hace el botón
``Aplicar'' de muchos reportes.

\subsection{\texorpdfstring{\codigo{req()}}{req()} y \texorpdfstring{\codigo{validate(need())}}{validate(need())}: condiciones antes de calcular}

Al abrir una app, algunas entradas tardan un instante en tener valor; otras
veces el usuario deja un campo vacío o elige filtros que no dejan ningún
dato. Dos herramientas evitan errores rojos en la pantalla:

\begin{itemize}
  \item \codigo{req(x)} detiene el cálculo \emph{en silencio} si \codigo{x}
        está vacío (\codigo{NULL}, \codigo{NA}, \codigo{""},
        \codigo{FALSE}). La salida simplemente se queda en blanco.
  \item \codigo{validate(need(condicion, "mensaje"))} detiene el cálculo y
        \emph{muestra el mensaje} en gris en lugar de la salida. Úsalo cuando
        el usuario necesite saber qué pasó.
\end{itemize}

\begin{rcode}
server_ticket <- function(input, output, session) {
  output$ticket <- renderText({
    req(input$metodo)                           # sin método: en blanco
    datos <- filter(transacciones, metodo_pago == input$metodo,
                    total_transaccion >= input$minimo)
    validate(need(nrow(datos) > 0,
                  "No hay transacciones con ese monto mínimo."))
    paste("Ticket promedio:", dollar(mean(datos$total_transaccion)))
  })
}

testServer(server_ticket, {
  session$setInputs(metodo = "Transferencia", minimo = 0)
  print(output$ticket)
  session$setInputs(minimo = 5000)              # nadie gasta tanto
  tryCatch(output$ticket, error = function(e) print(conditionMessage(e)))
})
\end{rcode}
\begin{rsalida}
[1] "Ticket promedio: $825.80"
[1] "No hay transacciones con ese monto mínimo."
\end{rsalida}

El ticket promedio con transferencia es de \$825.80. Con un mínimo de
\$5,000 no queda ninguna transacción y, en vez de un error críptico, el
usuario verá el mensaje que escribimos.

\subsection{\texorpdfstring{\codigo{isolate()}}{isolate()}: leer sin depender}

A veces quieres \emph{leer} una entrada sin que su cambio dispare un
recálculo. Por ejemplo, el título de una gráfica puede mencionar la meta
capturada, pero no quieres redibujar la gráfica cada vez que el usuario
teclea un dígito de la meta. \codigo{isolate(input\$meta)} lee el valor
actual sin crear la dependencia.

En la consola, fuera de una app, \codigo{isolate()} también sirve para
``asomarte'' a un valor reactivo:

\begin{rcode}
precio <- reactiveVal(500)                   # un valor reactivo
precio_con_iva <- reactive(precio() * 1.16)  # depende de precio
isolate(precio_con_iva())
precio(1000)                                 # cambiamos la fuente
isolate(precio_con_iva())                    # se recalcula solo
\end{rcode}
\begin{rsalida}
[1] 580
[1] 1160
\end{rsalida}

\subsection{Dos errores típicos de reactividad}
\label{sec:m6-errores-reactivos}

El primero: \textbf{llamar a un reactivo fuera de un contexto reactivo}.
Los reactivos solo pueden leerse dentro de un \codigo{render...()}, un
\codigo{reactive()}, un \codigo{observe...()} o con \codigo{isolate()}. Si lo
intentas en la consola o en el cuerpo del \codigo{server}, falla:

\begin{rcode}
tryCatch(precio_con_iva(),
         error = function(e) cat(conditionMessage(e), "\n"))
\end{rcode}
\begin{rsalida}
Operation not allowed without an active reactive context.
* You tried to do something that can only be done from inside a reactive consumer. 
\end{rsalida}

El segundo es su pariente: \textbf{usar \codigo{input\$...} directamente en
el cuerpo del \codigo{server}}, fuera de cualquier \codigo{render} o
\codigo{reactive}. \codigo{input} es un objeto \codigo{reactiveValues}, y
sus valores tampoco se pueden leer fuera de un consumidor:

\begin{rcode}
valores <- reactiveValues(n = 5)     # input funciona igual que esto
tryCatch(valores$n * 2,
         error = function(e) cat(conditionMessage(e), "\n"))
\end{rcode}
\begin{rsalida}
Can't access reactive value 'n' outside of reactive consumer.
i Do you need to wrap inside reactive() or observe()? 
\end{rsalida}

En una app real ese error aparece al abrirla, con la línea culpable, por
ejemplo \codigo{server <- function(input, output, session) \{ doble <-
input\$n * 2 ... \}}. La consola muestra
\codigo{Error in input\$n: Can't access reactive value 'n' outside of reactive
consumer} y la app se cierra. La solución es la que sugiere el propio
mensaje: envolver el cálculo en \codigo{reactive()} (\codigo{doble <-
reactive(input\$n * 2)}) y usarlo como \codigo{doble()}.

\begin{table}[H]
\centering
\caption{Herramientas de reactividad: cuál usar.}
\label{tab:m6-reactividad}
\small
\begin{tabularx}{\textwidth}{>{\raggedright\arraybackslash}p{3.6cm}
                             >{\raggedright\arraybackslash}X
                             >{\raggedright\arraybackslash}p{3.4cm}}
\toprule
\textbf{Función} & \textbf{Úsala cuando\ldots} & \textbf{Devuelve} \\
\midrule
\codigo{reactive()} & Varias salidas comparten un cálculo (datos
filtrados, KPIs) & Un valor que se lee con \codigo{()} \\
\codigo{eventReactive()} & El cálculo debe esperar a un botón & Un valor
que se lee con \codigo{()} \\
\codigo{observeEvent()} & Hay que \emph{hacer} algo al pulsar un botón
(limpiar filtros, guardar, avisar) & Nada (efecto secundario) \\
\codigo{observe()} & Hay que hacer algo cada vez que cambie cualquier cosa
que use & Nada \\
\codigo{reactiveVal()}, \codigo{reactiveValues()} & Necesitas un valor que
el propio servidor modifique (contadores, historial) & Un valor reactivo \\
\codigo{req()} & Falta un valor y la salida debe quedar en blanco & --- \\
\codigo{validate(need())} & Hay que explicar al usuario por qué no hay
resultado & --- \\
\codigo{isolate()} & Leer un valor sin crear dependencia & El valor \\
\bottomrule
\end{tabularx}
\end{table}

\begin{hazlotu}
Modifica \codigo{server\_ticket} para que, además del ticket promedio,
muestre cuántas transacciones cumplen la condición. Pruébalo con
\codigo{testServer()} usando \codigo{metodo = "Efectivo"} y \codigo{minimo
= 1000}.
\end{hazlotu}

% ----------------------------------------------------------------------------
\section{Organizar la pantalla: layouts}
\label{sec:m6-layouts}

\subsection{La cuadrícula de 12 columnas}

Shiny usa por debajo \termino{Bootstrap}, una biblioteca de diseño web que
divide el ancho de la pantalla en \textbf{12 columnas}. Con
\codigo{fluidRow()} creas una fila y con \codigo{column(ancho, ...)} repartes
esas 12 columnas: dos columnas de 6, tres de 4, una de 3 y otra de 9, etc.
(figura~\ref{fig:m6-cuadricula}). En pantallas angostas (un celular) las
columnas se apilan una debajo de otra automáticamente.

\begin{figure}[H]
\centering
\begin{tikzpicture}[font=\sffamily\footnotesize,
    col/.style={draw=white, line width=1.5pt, minimum height=0.7cm, align=center}]
  \foreach \i in {1,...,12} {
    \node[col, fill=gray!15, minimum width=1.2cm] at ({(\i-0.5)*1.2},3) {\i};
  }
  \node[col, fill=primario!25, minimum width=4.8cm] at (2.4,2) {\codigo{column(4, ...)}};
  \node[col, fill=secundario!25, minimum width=9.6cm] at (9.6,2) {\codigo{column(8, ...)}};
  \foreach \i in {0,...,3} {
    \node[col, fill=acento!25, minimum width=3.6cm] at ({1.8+\i*3.6},1) {\codigo{column(3, ...)}};
  }
  \node[col, fill=exito!20, minimum width=14.4cm] at (7.2,0) {\codigo{column(12, ...)}: todo el ancho};
  \node[anchor=east, text=gristexto] at (-0.1,3) {Bootstrap};
  \node[anchor=east, text=gristexto] at (-0.1,2) {Fila 1};
  \node[anchor=east, text=gristexto] at (-0.1,1) {Fila 2};
  \node[anchor=east, text=gristexto] at (-0.1,0) {Fila 3};
\end{tikzpicture}
\caption{Cada \codigo{fluidRow()} reparte 12 columnas entre sus
\codigo{column()}. Los anchos de una fila deben sumar 12 o menos.}
\label{fig:m6-cuadricula}
\end{figure}

\begin{rcode}
fluidRow(
  column(4, "Aquí van los filtros"),
  column(8, "Aquí va la gráfica")
)
\end{rcode}
\begin{rsalida}
<div class="row">
  <div class="col-sm-4">Aquí van los filtros</div>
  <div class="col-sm-8">Aquí va la gráfica</div>
</div>
\end{rsalida}

Cada \codigo{column(4, ...)} se convierte en un \codigo{<div
class="col-sm-4">}: la clase le dice a Bootstrap que ocupe 4 de las 12
columnas en pantallas pequeñas o mayores (\codigo{sm}, de \emph{small}).

\subsection{Los layouts que más usarás}

\begin{itemize}
  \item \codigo{fluidPage()}: la página base; ocupa todo el ancho.
  \item \codigo{sidebarLayout(sidebarPanel(...), mainPanel(...))}: la
        estructura clásica ``filtros a la izquierda, resultados a la
        derecha''. Por defecto la barra ocupa 4 columnas y el panel 8; con
        \codigo{width = 3} y \codigo{width = 9} das más espacio a las
        gráficas.
  \item \codigo{tabsetPanel(tabPanel("Título", ...), ...)}: pestañas
        dentro de una zona de la página.
  \item \codigo{navbarPage("Título", tabPanel(...), ...)}: una barra de
        navegación superior en la que cada \codigo{tabPanel()} es una
        ``página'' completa.
\end{itemize}

Este esqueleto combina todos (construir la \codigo{ui} es seguro en un
script; solo \codigo{shinyApp()} lanza la app):

\begin{rcode}
ui_esqueleto <- navbarPage(
  "Ventas 2023",
  tabPanel("Explorar",
    sidebarLayout(
      sidebarPanel(width = 3,
                   selectInput("tienda", "Tienda:", c("Todas", "Centro"))),
      mainPanel(width = 9,
        tabsetPanel(
          tabPanel("Gráfica", plotOutput("grafica")),
          tabPanel("Tabla", tableOutput("tabla"))
        )
      )
    )
  ),
  tabPanel("Acerca de", p("Datos simulados del curso."))
)
class(ui_esqueleto)
\end{rcode}
\begin{rsalida}
[1] "shiny.tag.list" "list"          
\end{rsalida}

La \codigo{ui} completa es una ``lista de etiquetas'' HTML
(\codigo{shiny.tag.list}). Nada se dibuja hasta que el navegador la recibe.

\begin{advertencia}
El error más frecuente al armar layouts es \textbf{olvidar una coma} entre
dos elementos, o dejar una coma \emph{de más} después del último. Una coma
faltante produce \codigo{Error: unexpected symbol} al ejecutar el archivo;
RStudio además marca la línea con una cruz roja. Revisa siempre la línea
anterior a la que señala el error. Sangrar el código (\tecla{Ctrl} +
\tecla{I} en RStudio) ayuda a ver qué está dentro de qué.
\end{advertencia}

% ----------------------------------------------------------------------------
\section{Caso práctico 1: el explorador de ventas}
\label{sec:m6-app02}

\begin{ejemplo}{Explorador de ventas 2023}
El equipo comercial quiere explorar las 1,000 transacciones de 2023 por
tienda, categoría y periodo; ver la evolución mensual medida en ventas,
unidades o número de transacciones; y consultar el top de productos. La app
completa está en \archivo{latex/codigo/apps/02-explorador-ventas/app.R}
(unas 160 líneas) y su distribución se muestra en la
figura~\ref{fig:m6-maqueta-app02}.
\end{ejemplo}

\begin{figure}[H]
\centering
\begin{tikzpicture}[font=\sffamily\footnotesize,
    zona/.style={draw=gristexto!60, rounded corners=2pt, align=center}]
  \draw[gristexto, rounded corners=3pt, line width=0.8pt] (0,0) rectangle (13,7.2);
  \node[anchor=west, font=\sffamily\large] at (0.2,6.8) {Explorador de ventas 2023};
  % sidebar
  \node[zona, fill=gray!10, minimum width=3.2cm, minimum height=6.1cm,
        anchor=north west] at (0.2,6.4) {};
  \foreach \y/\t in {6.0/Tienda: [Todas \ding{116}], 5.0/Categorías: [ Todas ],
                     4.0/Periodo: [01/01 a 31/12], 2.9/Medir por:\\ (o) Ventas\\ ( ) Unidades\\ ( ) Transacciones,
                     1.2/Productos en el top:\\ \textbf{---o---} 10} {
    \node[anchor=north west, align=left] at (0.35,\y) {\t};
  }
  % main
  \node[anchor=west, font=\sffamily\bfseries] at (3.7,6.1)
       {1,000 transacciones | Ventas: \$959,993};
  \draw[gristexto!60] (3.7,5.55) -- (12.8,5.55);
  \node[zona, fill=secundario!15, anchor=south west, minimum width=2.6cm,
        minimum height=0.45cm] at (3.7,5.55) {Evolución mensual};
  \node[zona, fill=white, anchor=south west, minimum width=2.3cm,
        minimum height=0.45cm] at (6.35,5.55) {Top productos};
  \node[zona, fill=secundario!8, minimum width=9.1cm, minimum height=5.0cm,
        anchor=north west] at (3.7,5.4) {};
  \foreach \x/\h in {0/2.6,1/2.1,2/2.9,3/1.8,4/3.0,5/2.8,6/2.7,7/2.5,8/2.2,9/2.2,10/2.6,11/3.1} {
    \fill[color=primario!60] ({4.1+\x*0.72},0.8) rectangle ++(0.5,\h);
  }
  \draw[gristexto] (3.9,0.8) -- (12.6,0.8);
  % callouts
  \node[text=acento, anchor=west] at (9.2,4.8) {\codigo{plotOutput("grafica")}};
  \node[text=acento, anchor=west, font=\sffamily\scriptsize] at (3.75,6.45) {\codigo{textOutput("resumen")}};
\end{tikzpicture}
\caption{Maqueta de la app 02: filtros en la barra lateral, resumen arriba y
resultados en dos pestañas.}
\label{fig:m6-maqueta-app02}
\end{figure}

Recorreremos el archivo en el mismo orden en que está escrito. Esa
organización (paquetes $\rightarrow$ datos $\rightarrow$ funciones
$\rightarrow$ interfaz $\rightarrow$ servidor) es la que te recomiendo para
cualquier app.

\subsection{Datos: se leen una sola vez, al arrancar}

Todo lo que está \emph{fuera} del \codigo{server} se ejecuta una sola vez,
cuando arranca la app, y lo comparten todos los usuarios. Por eso ahí leemos
y unimos los datos: si lo hiciéramos dentro del \codigo{server}, cada
persona que abra la app volvería a leer los CSV. Usamos
\codigo{buscar\_datasets()} (sección~\ref{sec:m6-rutas}) y
\codigo{floor\_date()} de \paquete{lubridate} para tener el primer día de
cada mes (ideal para graficar por mes con un eje de fechas).

\begin{rcode}
carpeta_datos <- buscar_datasets()
tiendas   <- read_csv(file.path(carpeta_datos, "tiendas.csv"),
                      show_col_types = FALSE)
productos <- read_csv(file.path(carpeta_datos, "productos.csv"),
                      show_col_types = FALSE)
ventas <- read_csv(file.path(carpeta_datos, "transacciones.csv"),
                   show_col_types = FALSE) %>%
  left_join(select(tiendas, tienda_id, nombre_tienda), by = "tienda_id") %>%
  left_join(select(productos, producto_id, nombre_producto, categoria),
            by = "producto_id") %>%
  mutate(mes = floor_date(fecha, "month"))   # primer día de cada mes

# Opciones para los filtros (se calculan una vez a partir de los datos)
opciones_tienda    <- c("Todas", sort(unique(ventas$nombre_tienda)))
opciones_categoria <- sort(unique(ventas$categoria))

opciones_tienda[1:4]
ventas %>%
  select(fecha, mes, nombre_tienda, categoria, total_transaccion) %>%
  head(3)
\end{rcode}
\begin{rsalida}
[1] "Todas"            "Sucursal Centro"  "Sucursal Este"   
[4] "Sucursal Express"
# A tibble: 3 x 5
  fecha      mes        nombre_tienda    categoria  total_transaccion
  <date>     <date>     <chr>            <chr>                  <dbl>
1 2023-01-13 2023-01-01 Sucursal Express Automotriz             1216.
2 2023-01-09 2023-01-01 Sucursal Norte   Automotriz              232.
3 2023-02-20 2023-02-01 Sucursal Plaza B Juguetes                826.
\end{rsalida}

Agregamos la opción \codigo{"Todas"} al principio del menú de tiendas: será
el valor por defecto y significará ``no filtrar''.

\subsection{Funciones de apoyo: la lógica, separada de la interfaz}

Esta es la decisión de diseño más importante del módulo: \textbf{todo el
trabajo con datos va en funciones normales de R}, que reciben datos y
parámetros y devuelven un data frame o una gráfica. No saben nada de
\codigo{input} ni de \codigo{output}. Ventajas: las pruebas en la consola
(sin abrir la app), las reutilizas en un reporte o en otra app, y el
\codigo{server} queda cortito y legible.

\begin{rcode}
filtrar_ventas <- function(datos, tienda = "Todas", categorias = NULL,
                           fechas = range(datos$fecha)) {
  resultado <- filter(datos, fecha >= fechas[1], fecha <= fechas[2])
  if (tienda != "Todas") {                 # "Todas" = no filtrar tienda
    resultado <- filter(resultado, nombre_tienda == tienda)
  }
  if (length(categorias) > 0) {            # vacío = todas las categorías
    resultado <- filter(resultado, categoria %in% categorias)
  }
  resultado
}

resumir_por_mes <- function(datos, metrica = "ventas") {
  datos %>%
    group_by(mes) %>%
    summarise(ventas        = sum(total_transaccion),
              unidades      = sum(cantidad),
              transacciones = n(),
              .groups = "drop") %>%
    mutate(valor = .data[[metrica]])   # la columna elegida por el usuario
}
\end{rcode}

\codigo{resumir\_por\_mes()} calcula las tres métricas y copia en la columna
\codigo{valor} la que eligió el usuario. \codigo{.data[[metrica]]} es la
forma de \paquete{dplyr} de usar una columna cuyo \emph{nombre} está
guardado en una variable de texto: justo lo que llega de un
\codigo{radioButtons()}. Probemos ambas funciones como lo haría la app con
``Sucursal Norte'' y dos categorías:

\begin{rcode}
norte <- filtrar_ventas(ventas, tienda = "Sucursal Norte",
                        categorias = c("Electrónica", "Hogar"))
nrow(norte)
resumir_por_mes(norte, metrica = "unidades") %>% head(4)
\end{rcode}
\begin{rsalida}
[1] 21
# A tibble: 4 x 5
  mes        ventas unidades transacciones valor
  <date>      <dbl>    <dbl>         <int> <dbl>
1 2023-01-01  2741.        5             1     5
2 2023-02-01   871.        2             2     2
3 2023-03-01  1728.        4             2     4
4 2023-04-01   153.        1             1     1
\end{rsalida}

\begin{rcode}
grafica_mensual <- function(resumen, metrica = "ventas") {
  etiquetas <- c(ventas = "Ventas", unidades = "Unidades vendidas",
                 transacciones = "Transacciones")
  formato <- if (metrica == "ventas") label_dollar() else label_comma()
  ggplot(resumen, aes(x = mes, y = valor)) +
    geom_col(fill = "#2a78d6", width = 22) +    # ancho en días
    scale_x_date(date_labels = "%b", date_breaks = "1 month") +
    scale_y_continuous(labels = formato,
                       expand = expansion(mult = c(0, 0.05))) +
    labs(x = NULL, y = etiquetas[[metrica]]) +
    theme_minimal(base_size = 13) +
    theme(panel.grid.minor = element_blank(),
          panel.grid.major.x = element_blank())
}

top_productos <- function(datos, n = 10) {
  datos %>%
    group_by(Producto = nombre_producto, `Categoría` = categoria) %>%
    summarise(Unidades = sum(cantidad),
              Ventas   = sum(total_transaccion), .groups = "drop") %>%
    slice_max(Ventas, n = n, with_ties = FALSE) %>%
    mutate(Ventas = dollar(Ventas))       # texto con formato de moneda
}
\end{rcode}

Ahora la gráfica que verá el usuario al abrir la app, con todos los datos y
la métrica por defecto (figura~\ref{fig:m06-explorador-mensual}):

\begin{rcode}
p_mensual <- resumir_por_mes(ventas, "ventas") %>%
  grafica_mensual("ventas")
p_mensual
\end{rcode}

\begin{figure}[H]
\centering
\includegraphics[width=0.9\textwidth]{m06-explorador-mensual}
\caption{Pestaña ``Evolución mensual'' de la app 02 con los filtros por
defecto (salida de \codigo{grafica\_mensual()}).}
\label{fig:m06-explorador-mensual}
\end{figure}

Una sola serie, un solo color, barras que empiezan en cero y eje con
formato de moneda. Se ve que abril fue el mes más flojo y diciembre el más
fuerte. La tabla del top de productos, con el valor por defecto del
deslizador (aquí pedimos 5 para ahorrar espacio):

\begin{rcode}
top_productos(ventas, n = 5)
\end{rcode}
\begin{rsalida}
# A tibble: 5 x 4
  Producto    Categoría  Unidades Ventas    
  <chr>       <chr>         <dbl> <chr>     
1 Producto 06 Juguetes         82 $34,028.04
2 Producto 19 Deportes         85 $28,367.95
3 Producto 31 Alimentos        74 $25,429.70
4 Producto 47 Automotriz       82 $24,527.21
5 Producto 46 Automotriz       74 $24,470.29
\end{rsalida}

El Producto~06 (Juguetes) lidera con \$34,028. Convertimos \codigo{Ventas}
a texto con \codigo{dollar()} al final porque \codigo{renderTable()} muestra
los números tal cual; así el usuario ve ``\$34,028.04'' en vez de
``34028.04''.

\subsection{La interfaz}

Con la lógica lista, la \codigo{ui} solo acomoda piezas. Observa que todas
las opciones salen de los datos (\codigo{opciones\_tienda},
\codigo{min(ventas\$fecha)}\ldots) y que el \codigo{selectizeInput()} usa
\codigo{placeholder = "Todas"} para que el cuadro vacío se lea como ``todas
las categorías''.

\begin{rnoejecutar}[title=app.R (parte 5): interfaz]
ui <- fluidPage(
  titlePanel("Explorador de ventas 2023"),
  sidebarLayout(
    # Barra lateral: todos los filtros juntos
    sidebarPanel(
      width = 3,
      selectInput("tienda", "Tienda:", choices = opciones_tienda),
      selectizeInput("categorias", "Categorías:",
                     choices = opciones_categoria, multiple = TRUE,
                     options = list(placeholder = "Todas")),
      dateRangeInput("fechas", "Periodo:",
                     start = min(ventas$fecha), end = max(ventas$fecha),
                     min = min(ventas$fecha), max = max(ventas$fecha),
                     format = "dd/mm/yyyy", language = "es",
                     separator = " a "),
      radioButtons("metrica", "Medir por:",
                   choices = c("Ventas ($)" = "ventas",
                               "Unidades" = "unidades",
                               "Transacciones" = "transacciones")),
      sliderInput("n_top", "Productos en el top:",
                  min = 5, max = 20, value = 10, step = 5)
    ),
    # Panel principal: resumen arriba y resultados en pestañas
    mainPanel(
      width = 9,
      h4(textOutput("resumen")),
      tabsetPanel(
        tabPanel("Evolución mensual",
                 plotOutput("grafica", height = "380px")),
        tabPanel("Top productos", tableOutput("tabla_top"))
      )
    )
  )
)
\end{rnoejecutar}

\subsection{El servidor}

Gracias a las funciones, el \codigo{server} cabe en una pantalla. Tiene un
único reactivo, \codigo{datos\_filtrados()}, que alimenta las tres salidas
(el grafo de la figura~\ref{fig:m6-grafo}):

\begin{rnoejecutar}[title=app.R (partes 6 y 7): servidor y lanzamiento]
server <- function(input, output, session) {

  # Expresión reactiva: filtra UNA vez y la reutilizan las tres salidas
  datos_filtrados <- reactive({
    req(input$fechas)                      # espera a que haya fechas
    filtrar_ventas(ventas, input$tienda, input$categorias, input$fechas)
  })

  output$resumen <- renderText({
    d <- datos_filtrados()                 # ¡con paréntesis!
    paste0(comma(nrow(d)), " transacciones  |  Ventas: ",
           dollar(sum(d$total_transaccion), accuracy = 1))
  })

  output$grafica <- renderPlot({
    validate(need(nrow(datos_filtrados()) > 0,
                  "No hay ventas con estos filtros. Amplía el periodo."))
    resumir_por_mes(datos_filtrados(), input$metrica) %>%
      grafica_mensual(input$metrica)
  }, res = 96)

  output$tabla_top <- renderTable({
    top_productos(datos_filtrados(), n = input$n_top)
  }, align = "llrr", digits = 0)
}

# 7. Lanzar la app -------------------------------------------------------
shinyApp(ui, server)
\end{rnoejecutar}

Detalles que vale la pena notar:

\begin{itemize}
  \item \codigo{req(input\$fechas)} evita calcular en el instante en que la
        app arranca y el calendario aún no envía su valor.
  \item \codigo{validate(need(...))} muestra un mensaje amable si el
        periodo y los filtros no dejan ninguna venta.
  \item \codigo{renderPlot(\{...\}, res = 96)} dibuja la gráfica con una
        resolución adecuada para pantalla (textos de tamaño normal).
  \item \codigo{renderTable(..., align = "llrr", digits = 0)} alinea a la
        izquierda las dos columnas de texto y a la derecha las numéricas, y
        muestra las unidades sin decimales.
\end{itemize}

\subsection{Probar la app completa con \texorpdfstring{\codigo{testServer()}}{testServer()}}

\codigo{testServer()} también acepta una app completa: le pasas
\codigo{shinyAppDir("carpeta")} y ejecuta su \codigo{app.R} (incluida la
lectura de datos) sin abrir el navegador. Dentro del bloque puedes leer
incluso los reactivos internos, como \codigo{datos\_filtrados()}:

\begin{rcode}
app02 <- shinyAppDir("latex/codigo/apps/02-explorador-ventas")
testServer(app02, {
  session$setInputs(tienda = "Todas", categorias = NULL,
                    fechas = as.Date(c("2023-01-01", "2023-12-31")),
                    metrica = "ventas", n_top = 10)
  print(output$resumen)
  session$setInputs(tienda = "Sucursal Centro",
                    fechas = as.Date(c("2023-07-01", "2023-12-31")))
  print(output$resumen)
  print(range(datos_filtrados()$fecha))       # un reactivo interno
})
\end{rcode}
\begin{rsalida}
[1] "1,000 transacciones  |  Ventas: $959,993"
[1] "46 transacciones  |  Ventas: $49,559"
[1] "2023-07-05" "2023-12-31"
\end{rsalida}

Con todos los filtros por defecto la app resume las 1,000 transacciones
(\$959,993); al pasar a la Sucursal Centro en el segundo semestre quedan
46 transacciones por \$49,559, y el rango de fechas (del 5 de julio al 31
de diciembre) confirma que el filtro funciona. Si algo fallara (una columna mal escrita, un paquete faltante),
lo verías aquí como error, antes de que lo vea tu directora.

\begin{hazlotu}
Abre la app con \codigo{shiny::runApp("latex/codigo/apps/02-explorador-ventas")},
elige ``Sucursal Outlet'', la categoría ``Belleza'' y la métrica
``Unidades''. ¿En qué mes vendió más unidades de belleza el Outlet? Luego
reduce el periodo a una semana de enero y observa el mensaje de
\codigo{validate()}.
\end{hazlotu}

% ----------------------------------------------------------------------------
\section{Dashboards profesionales con shinydashboard}
\label{sec:m6-shinydashboard}

\subsection{Las piezas de \texorpdfstring{\paquete{shinydashboard}}{shinydashboard}}

\paquete{shinydashboard} es un paquete que da a tu app el aspecto clásico
de un tablero de BI: encabezado de color, menú lateral y cajas. Su
estructura siempre es la misma:

\begin{itemize}
  \item \codigo{dashboardPage(header, sidebar, body)}: la página completa,
        con tres argumentos obligatorios en ese orden.
  \item \codigo{dashboardHeader(title = ...)}: la barra superior.
  \item \codigo{dashboardSidebar(...)}: la barra lateral. Suele contener un
        \codigo{sidebarMenu()} con varios \codigo{menuItem()} (las
        ``pestañas'') y, debajo, los filtros globales.
  \item \codigo{dashboardBody(...)}: el cuerpo. Con
        \codigo{tabItems(tabItem(tabName = ...), ...)} defines el contenido
        de cada pestaña; el \codigo{tabName} debe coincidir con el del
        \codigo{menuItem()}.
  \item \codigo{box(title, width, status, ...)}: una caja con título que
        agrupa una gráfica o tabla. \codigo{width} usa la misma cuadrícula de
        12 columnas; \codigo{status = "primary"} le da el color del tema.
  \item \codigo{valueBox(valor, subtitulo, icon, color)}: la tarjeta grande
        de KPI. \codigo{infoBox()} es una variante más compacta con el ícono
        a la izquierda.
\end{itemize}

Un \codigo{valueBox()} también es solo HTML. Imprímelo para ver su
estructura: el valor va en un encabezado grande (\codigo{h3}) y el
subtítulo en un párrafo.

\begin{rcode}
valueBox(value = "$959,993", subtitle = "Ventas totales",
         icon = icon("dollar-sign"), color = "blue", width = 3)
\end{rcode}
\begin{rsalida}
<div class="col-sm-3">
  <div class="small-box bg-blue">
    <div class="inner">
      <h3>$959,993</h3>
      <p>Ventas totales</p>
    </div>
    <div class="icon-large">
      <i class="fas fa-dollar-sign" role="presentation" aria-label="dollar-sign icon"></i>
    </div>
  </div>
</div>
\end{rsalida}

Los íconos vienen de la colección gratuita Font Awesome: busca nombres en
\url{https://fontawesome.com/icons} (por ejemplo \codigo{"store"},
\codigo{"users"}, \codigo{"chart-line"}) y úsalos con \codigo{icon()}.

\begin{consejo}
Resiste la tentación de pintar cada \codigo{valueBox()} de un color
distinto (verde, rojo, amarillo, morado\ldots). En un dashboard el color
debe \emph{significar} algo: usa un mismo color para todas las tarjetas y
reserva el rojo o el verde para cuando un KPI esté por debajo o por encima
de su meta (lo harás en el ejercicio~\ref{ej:m6-3}).
\end{consejo}

% ----------------------------------------------------------------------------
\section{Caso práctico 2: el dashboard ejecutivo}
\label{sec:m6-app03}

\begin{ejemplo}{Dashboard ejecutivo de ventas}
La dirección quiere un tablero con cinco pestañas: \textbf{Resumen} (4 KPIs
y evolución mensual), \textbf{Productos}, \textbf{Tiendas},
\textbf{Clientes} y \textbf{Datos} (tabla con búsqueda y descarga). Los
filtros de periodo, tiendas y categorías deben afectar a todas las pestañas,
y un botón debe regresar todo a su estado inicial. La app está en
\archivo{latex/codigo/apps/03-dashboard-ejecutivo/} y su distribución
aparece en la figura~\ref{fig:m6-maqueta-app03}.
\end{ejemplo}

\begin{figure}[H]
\centering
\begin{tikzpicture}[font=\sffamily\footnotesize,
    zona/.style={draw=gristexto!50, rounded corners=1.5pt, align=center},
    kpi/.style={zona, fill=primario!70, text=white, minimum width=2.25cm,
                minimum height=1.05cm}]
  \draw[gristexto, line width=0.8pt] (0,0) rectangle (13.4,7.6);
  \fill[primario] (0,7) rectangle (3.2,7.6);
  \fill[primario!80] (3.2,7) rectangle (13.4,7.6);
  \node[white, font=\sffamily\bfseries] at (1.6,7.3) {Ventas 2023};
  \node[white] at (3.6,7.3) {$\equiv$};
  % sidebar
  \fill[gray!75!black] (0,0) rectangle (3.2,7);
  \foreach \y/\t in {6.6/Resumen, 6.15/Productos, 5.7/Tiendas, 5.25/Clientes, 4.8/Datos} {
    \node[anchor=west, text=white] at (0.25,\y) {\t};
  }
  \fill[primario!60, opacity=0.6] (0,6.4) rectangle (3.2,6.8);
  \node[anchor=west, text=white, align=left] at (0.2,3.95) {Periodo:\\[1pt]
    \colorbox{white}{\textcolor{black}{01/01/2023 a 31/12/2023}}};
  \node[anchor=west, text=white, align=left] at (0.2,2.95) {Tiendas:\\[1pt]
    \colorbox{white}{\textcolor{black}{\makebox[2.3cm][l]{Todas}}}};
  \node[anchor=west, text=white, align=left] at (0.2,1.95) {Categorías:\\[1pt]
    \colorbox{white}{\textcolor{black}{\makebox[2.3cm][l]{Todas}}}};
  \node[anchor=west, zona, fill=white] at (0.25,1.0) {Quitar filtros};
  % body
  \fill[gray!8] (3.2,0) rectangle (13.4,7);
  \foreach \i/\v/\t in {0/\$959,993/Ventas totales, 1/1{,}000/Transacciones,
                        2/\$960/Ticket promedio, 3/198/Clientes activos} {
    \node[kpi, anchor=north west, align=left] at ({3.4+\i*2.5},6.8)
      {\textbf{\normalsize\v}\\ \t};
  }
  \node[zona, fill=white, minimum width=6.4cm, minimum height=4.2cm,
        anchor=north west] (evo) at (3.4,5.45) {};
  \fill[primario] (3.4,5.45) rectangle (9.8,5.0);
  \node[white, anchor=west] at (3.5,5.22) {Ventas por mes};
  \draw[primario!70, thick] (3.8,3.0) -- (4.3,2.5) -- (4.8,3.2) -- (5.3,2.0)
       -- (5.8,3.3) -- (6.3,3.1) -- (6.8,3.0) -- (7.3,2.8) -- (7.8,2.4)
       -- (8.3,2.45) -- (8.8,2.9) -- (9.3,3.45);
  \fill[acento] (9.3,3.45) circle (0.08);
  \node[zona, fill=white, minimum width=3.4cm, minimum height=4.2cm,
        anchor=north west] at (9.85,5.45) {};
  \fill[primario] (9.85,5.45) rectangle (13.25,5.0);
  \node[white, anchor=west, font=\sffamily\scriptsize] at (9.9,5.22) {Ventas por método de pago};
  \foreach \y/\w in {4.4/2.6, 3.8/2.2, 3.2/2.0, 2.6/0.9} {
    \fill[primario!70] (10.2,\y) rectangle ++(\w,0.35);
  }
  \node[text=gristexto] at (8.3,0.6) {\codigo{box(width = 8)} \hspace{1.6cm} \codigo{box(width = 4)}};
  \node[text=acento, anchor=south] at (8.3,6.85) {};
\end{tikzpicture}
\caption{Maqueta de la pestaña ``Resumen'' de la app 03: menú y filtros en
la barra lateral, cuatro \codigo{valueBox()} arriba y dos \codigo{box()} con
gráficas.}
\label{fig:m6-maqueta-app03}
\end{figure}

La app tiene dos archivos:

\begin{textocodigo}
03-dashboard-ejecutivo/
|-- app.R                      <- interfaz y servidor
+-- R/
    +-- funciones_dashboard.R  <- datos, KPIs, gráficas y tablas
\end{textocodigo}

Shiny tiene una regla muy útil: \textbf{antes de ejecutar \archivo{app.R},
carga automáticamente todos los archivos \archivo{.R} de la subcarpeta
\archivo{R/}}. Así la lógica vive en su propio archivo y \archivo{app.R} se
queda con lo visual. Vamos a recorrer primero las funciones (que sí podemos
ejecutar y probar aquí) y después la app.

\subsection{Cargar y unir los datos}

\codigo{cargar\_datos\_dashboard()} lee cuatro CSV, se queda con las
columnas necesarias y los une a las transacciones con los \emph{joins} del
Módulo~\ref{cap:manipulacion}. El resultado es una sola tabla ``ancha'' con
todo lo que el dashboard necesita:

\begin{rcode}
# ---- Colores y tema corporativos ---------------------------------------
color_principal <- "#2a78d6"   # una sola serie
color_resalte   <- "#eb6834"   # para destacar un elemento
color_gris      <- "#b5b3ad"   # el resto, cuando se resalta uno
paleta_libro <- c("#2a78d6", "#eb6834", "#1baf7a", "#eda100",
                  "#e87ba4", "#008300", "#4a3aa7", "#e34948")

tema_dashboard <- function(base_size = 12) {
  theme_minimal(base_size = base_size) +
    theme(panel.grid.minor = element_blank(),
          legend.position  = "top",
          plot.margin      = margin(5, 15, 5, 5))  # aire a la derecha
}

cargar_datos_dashboard <- function(carpeta = buscar_datasets()) {
  leer <- function(archivo) {
    read_csv(file.path(carpeta, archivo), show_col_types = FALSE)
  }
  productos <- leer("productos.csv") %>%
    select(producto_id, nombre_producto, categoria, costo)
  tiendas <- leer("tiendas.csv") %>%
    transmute(tienda_id, ciudad,
              tienda = sub("Sucursal ", "", nombre_tienda))  # nombre corto
  clientes <- leer("clientes.csv") %>%
    select(cliente_id, nombre_cliente = nombre, segmento, es_premium)

  leer("transacciones.csv") %>%
    left_join(productos, by = "producto_id") %>%
    left_join(tiendas,   by = "tienda_id") %>%
    left_join(clientes,  by = "cliente_id") %>%
    mutate(mes = floor_date(fecha, "month"),
           segmento = factor(segmento, levels = c("Joven", "Adulto",
                                                  "Maduro", "Senior"))) %>%
    select(transaccion_id, fecha, mes, tienda, ciudad, categoria,
           producto_id, nombre_producto, cliente_id, nombre_cliente,
           segmento, es_premium, metodo_pago, cantidad, precio_venta,
           costo, total_transaccion)
}

datos <- cargar_datos_dashboard()
glimpse(datos)
\end{rcode}
\begin{rsalida}
Rows: 1,000
Columns: 17
$ transaccion_id    <dbl> 1, 2, 3, 4, 5, 6, 7, 8, 9, 10, 11, 12, 13,~
$ fecha             <date> 2023-01-13, 2023-01-09, 2023-02-20, 2023-~
$ mes               <date> 2023-01-01, 2023-01-01, 2023-02-01, 2023-~
$ tienda            <chr> "Express", "Norte", "Plaza B", "Plaza A", ~
$ ciudad            <chr> "Querétaro", "Guadalajara", "Guadalajara",~
$ categoria         <chr> "Automotriz", "Automotriz", "Juguetes", "D~
$ producto_id       <dbl> 3, 47, 6, 17, 28, 46, 30, 48, 21, 12, 49, ~
$ nombre_producto   <chr> "Producto 03", "Producto 47", "Producto 06~
$ cliente_id        <dbl> 155, 155, 75, 51, 173, 135, 90, 168, 13, 9~
$ nombre_cliente    <chr> "Marta González", "Marta González", "Diego~
$ segmento          <fct> Adulto, Adulto, Maduro, Senior, Maduro, Ma~
$ es_premium        <lgl> TRUE, TRUE, FALSE, FALSE, FALSE, TRUE, FAL~
$ metodo_pago       <chr> "Tarjeta Crédito", "Tarjeta Débito", "Tarj~
$ cantidad          <dbl> 4, 4, 3, 4, 1, 2, 4, 2, 2, 5, 2, 3, 4, 2, ~
$ precio_venta      <dbl> 303.97, 58.07, 275.46, 105.28, 205.47, 534~
$ costo             <dbl> 294.30, 55.00, 84.55, 156.61, 17.23, 177.5~
$ total_transaccion <dbl> 1215.88, 232.28, 826.38, 421.12, 205.47, 1~
\end{rsalida}

1,000 filas y 17 columnas: cada transacción ya trae el nombre corto de su
tienda, la categoría y el costo del producto, y el segmento del cliente.
(Usamos \codigo{sub("Sucursal ", "", nombre\_tienda)} para mostrar
``Centro'' en lugar de ``Sucursal Centro'' y ahorrar espacio en las
gráficas.) La función \codigo{buscar\_datasets()} de este archivo es la
misma de la sección~\ref{sec:m6-rutas}.

\subsection{Filtros globales y KPIs}

\begin{rcode}
# ---- Filtros globales --------------------------------------------------
# Vacío (NULL o character(0)) significa "todas".
filtrar_datos <- function(datos, fechas, tiendas = NULL,
                          categorias = NULL) {
  datos <- filter(datos, fecha >= fechas[1], fecha <= fechas[2])
  if (length(tiendas) > 0) {
    datos <- filter(datos, tienda %in% tiendas)
  }
  if (length(categorias) > 0) {
    datos <- filter(datos, categoria %in% categorias)
  }
  datos
}

calcular_kpis <- function(datos) {
  datos %>%
    summarise(
      ventas          = sum(total_transaccion),
      transacciones   = n(),
      ticket_promedio = ventas / transacciones,
      unidades        = sum(cantidad),
      clientes        = n_distinct(cliente_id),
      pct_premium     = mean(es_premium[!duplicated(cliente_id)]),
      compras_cliente = transacciones / clientes
    )
}

# Formatos para mostrar los KPIs como texto ("--" si no hay dato)
formato_pesos <- function(x) {
  if_else(is.na(x), "--", dollar(x, accuracy = 1))
}
formato_numero <- function(x, decimales = 0) {
  if_else(is.na(x), "--", comma(x, accuracy = 10^-decimales))
}
formato_pct <- function(x) {
  if_else(is.na(x), "--", percent(x, accuracy = 0.1))
}
\end{rcode}

Dos detalles de negocio en \codigo{calcular\_kpis()}: los \emph{clientes
activos} se cuentan con \codigo{n\_distinct()} (un cliente que compró cinco
veces cuenta una vez), y el \% de clientes premium se calcula sobre
clientes únicos (\codigo{!duplicated(cliente\_id)}), no sobre
transacciones. Las funciones de formato devuelven \codigo{"-{}-"} cuando no
hay dato (por ejemplo, el ticket promedio de un filtro sin ventas), en lugar
de un feo \codigo{NA}.

\begin{rcode}
anio <- as.Date(c("2023-01-01", "2023-12-31"))
kpis_total <- calcular_kpis(filtrar_datos(datos, anio))
kpis_total %>% select(ventas, transacciones, ticket_promedio, clientes)
formato_pesos(kpis_total$ventas)
formato_pct(kpis_total$pct_premium)

# Los mismos KPIs para dos tiendas y una categoría, primer semestre
filtrar_datos(datos, as.Date(c("2023-01-01", "2023-06-30")),
              tiendas = c("Centro", "Norte"),
              categorias = "Electrónica") %>%
  calcular_kpis() %>%
  select(ventas, transacciones, ticket_promedio, clientes)
\end{rcode}
\begin{rsalida}
# A tibble: 1 x 4
   ventas transacciones ticket_promedio clientes
    <dbl>         <int>           <dbl>    <int>
1 959993.          1000            960.      198
[1] "$959,993"
[1] "36.4%"
# A tibble: 1 x 4
  ventas transacciones ticket_promedio clientes
   <dbl>         <int>           <dbl>    <int>
1  6937.            10            694.       10
\end{rsalida}

En todo 2023 hubo 1,000 transacciones por \$959,993, con un ticket
promedio de \$960 y 198 de los 200 clientes activos; el 36.4\% de esos
clientes es premium. Con los filtros del segundo ejemplo la historia cambia
por completo: solo 10 transacciones de Electrónica entre Centro y Norte en el
primer semestre. Estas cifras son exactamente las que mostrarán los
\codigo{valueBox()}.

\subsection{Las gráficas de cada pestaña}

Cada gráfica es una función que recibe los datos \emph{ya filtrados} y
devuelve un objeto \codigo{ggplot}. Todas usan el tema del dashboard y los
colores corporativos del Módulo~\ref{cap:visualizacion}.

\begin{rcode}
grafica_evolucion <- function(datos) {
  por_mes <- datos %>%
    group_by(mes) %>%
    summarise(ventas = sum(total_transaccion), .groups = "drop")
  mejor_mes <- slice_max(por_mes, ventas, n = 1, with_ties = FALSE)
  ggplot(por_mes, aes(x = mes, y = ventas)) +
    geom_line(color = color_principal, linewidth = 0.8) +
    geom_point(color = color_principal, size = 2) +
    geom_point(data = mejor_mes, color = color_resalte, size = 3.5) +
    geom_text(data = mejor_mes, aes(label = dollar(ventas, accuracy = 1)),
              vjust = -1.2, size = 3.5, color = "gray25") +
    scale_x_date(date_labels = "%b", date_breaks = "1 month") +
    scale_y_continuous(labels = label_dollar(), limits = c(0, NA),
                       expand = expansion(mult = c(0, 0.15))) +
    labs(x = NULL, y = "Ventas") +
    tema_dashboard()
}

p_evolucion <- grafica_evolucion(datos)
p_evolucion
\end{rcode}

\begin{figure}[H]
\centering
\includegraphics[width=0.9\textwidth]{m06-evolucion}
\caption{Pestaña Resumen: ventas por mes con el mejor mes resaltado
(\codigo{grafica\_evolucion()}).}
\label{fig:m06-evolucion}
\end{figure}

La línea usa un solo color y solo el mejor mes lleva color de resalte y
etiqueta (diciembre, con \$96,394): el ojo va directo a lo importante. Como
el eje Y empieza en cero, la caída de abril se ve en su verdadera magnitud.

\begin{rcode}
grafica_metodo_pago <- function(datos) {
  datos %>%
    group_by(metodo_pago) %>%
    summarise(ventas = sum(total_transaccion), .groups = "drop") %>%
    mutate(participacion = ventas / sum(ventas)) %>%
    ggplot(aes(x = participacion,
               y = reorder(metodo_pago, participacion))) +
    geom_col(fill = color_principal, width = 0.7) +
    scale_x_continuous(labels = label_percent(),
                       expand = expansion(mult = c(0, 0.05))) +
    labs(x = "% de las ventas", y = NULL) +
    tema_dashboard()
}

grafica_barras <- function(resumen, etiqueta_y = NULL) {
  # resumen: columnas 'nombre' y 'ventas'; barras horizontales ordenadas
  ggplot(resumen, aes(x = ventas, y = reorder(nombre, ventas))) +
    geom_col(fill = color_principal, width = 0.7) +
    scale_x_continuous(labels = label_dollar(),
                       expand = expansion(mult = c(0, 0.05))) +
    labs(x = "Ventas", y = etiqueta_y) +
    tema_dashboard()
}

grafica_top_productos <- function(datos, n = 10) {
  datos %>%
    group_by(nombre = nombre_producto) %>%
    summarise(ventas = sum(total_transaccion), .groups = "drop") %>%
    slice_max(ventas, n = n, with_ties = FALSE) %>%
    grafica_barras()
}

grafica_categorias <- function(datos) {
  datos %>%
    group_by(nombre = categoria) %>%
    summarise(ventas = sum(total_transaccion), .groups = "drop") %>%
    grafica_barras()
}

p_top <- grafica_top_productos(datos)
p_top
\end{rcode}

\codigo{grafica\_barras()} es una función ``genérica'' que reutilizan la
gráfica de productos y la de categorías: ambas resumen a dos columnas,
\codigo{nombre} y \codigo{ventas}, y le pasan el resultado. Escribir la
gráfica una sola vez garantiza que las dos pestañas se vean iguales
(figura~\ref{fig:m06-top-productos}).

\begin{figure}[H]
\centering
\includegraphics[width=0.9\textwidth]{m06-top-productos}
\caption{Pestaña Productos: los 10 productos con más ventas
(\codigo{grafica\_top\_productos()}).}
\label{fig:m06-top-productos}
\end{figure}

\begin{rcode}
grafica_tiendas <- function(datos) {
  por_tienda <- datos %>%
    group_by(tienda) %>%
    summarise(ventas = sum(total_transaccion), .groups = "drop") %>%
    mutate(es_mejor = ventas == max(ventas))    # la tienda líder
  ggplot(por_tienda, aes(x = ventas, y = reorder(tienda, ventas),
                         fill = es_mejor)) +
    geom_col(width = 0.7) +
    geom_text(data = filter(por_tienda, es_mejor),
              aes(label = dollar(ventas, accuracy = 1)),
              hjust = 1.1, color = "white", size = 3.5) +
    scale_fill_manual(values = c(`TRUE` = color_resalte,
                                 `FALSE` = color_gris), guide = "none") +
    scale_x_continuous(labels = label_dollar(),
                       expand = expansion(mult = c(0, 0.05))) +
    labs(x = "Ventas", y = NULL) +
    tema_dashboard()
}

p_tiendas <- grafica_tiendas(datos)
p_tiendas
\end{rcode}

\begin{figure}[H]
\centering
\includegraphics[width=0.9\textwidth]{m06-tiendas}
\caption{Pestaña Tiendas: la tienda líder en color de resalte y el resto en
gris (\codigo{grafica\_tiendas()}).}
\label{fig:m06-tiendas}
\end{figure}

La figura~\ref{fig:m06-tiendas} aplica la técnica de ``resaltar uno y
agrisar el resto'': la pregunta es ``¿quién lidera?'' y la respuesta salta a
la vista. Un hallazgo interesante: el \emph{Outlet}, una de las tiendas más
pequeñas (350 m$^2$), es la que más vende en estos datos simulados, con
\$114,377.

\begin{rcode}
grafica_segmentos <- function(datos) {
  datos %>%
    mutate(tipo = if_else(es_premium, "Premium", "Regular")) %>%
    group_by(segmento, tipo) %>%
    summarise(ventas = sum(total_transaccion), .groups = "drop") %>%
    ggplot(aes(x = segmento, y = ventas, fill = tipo)) +
    geom_col(position = position_dodge(width = 0.8), width = 0.7) +
    scale_fill_manual(values = c(Premium = paleta_libro[2],
                                 Regular = paleta_libro[1])) +
    scale_y_continuous(labels = label_dollar(),
                       expand = expansion(mult = c(0, 0.05))) +
    labs(x = "Segmento de edad", y = "Ventas", fill = "Tipo de cliente") +
    tema_dashboard()
}

p_segmentos <- grafica_segmentos(datos)
p_segmentos
\end{rcode}

\begin{figure}[H]
\centering
\includegraphics[width=0.9\textwidth]{m06-segmentos}
\caption{Pestaña Clientes: ventas por segmento de edad y tipo de cliente
(\codigo{grafica\_segmentos()}).}
\label{fig:m06-segmentos}
\end{figure}

Aquí sí hay dos series (Premium y Regular), así que usamos dos colores de la
paleta y leyenda. En todos los segmentos los clientes regulares aportan más
ventas que los premium, lo cual es lógico porque son cerca del 64\% de los
clientes.

\subsection{Tablas y detalle}

\begin{rcode}
# ---- Tablas ------------------------------------------------------------
tabla_tiendas <- function(datos) {
  datos %>%
    group_by(Tienda = tienda) %>%
    summarise(Ventas = sum(total_transaccion),
              Transacciones = n(),
              Ticket = Ventas / Transacciones, .groups = "drop") %>%
    arrange(desc(Ventas)) %>%
    mutate(Ventas = formato_pesos(Ventas), Ticket = formato_pesos(Ticket))
}

tabla_top_clientes <- function(datos, n = 10) {
  datos %>%
    group_by(Cliente = nombre_cliente, Segmento = segmento) %>%
    summarise(Compras = n(), Ventas = sum(total_transaccion),
              .groups = "drop") %>%
    slice_max(Ventas, n = n, with_ties = FALSE) %>%
    mutate(Ventas = formato_pesos(Ventas))
}

# Datos listos para la tabla interactiva y para las descargas
datos_detalle <- function(datos) {
  datos %>%
    arrange(fecha, transaccion_id) %>%
    select(Fecha = fecha, Tienda = tienda, `Categoría` = categoria,
           Producto = nombre_producto, Cliente = nombre_cliente,
           Cantidad = cantidad, Precio = precio_venta,
           Total = total_transaccion, `Método de pago` = metodo_pago)
}

tabla_tiendas(datos) %>% head(4)
tabla_top_clientes(datos, n = 3)
\end{rcode}
\begin{rsalida}
# A tibble: 4 x 4
  Tienda  Ventas   Transacciones Ticket
  <chr>   <chr>            <int> <chr> 
1 Outlet  $114,377           106 $1,079
2 Sur     $105,305           109 $966  
3 Plaza A $105,216           109 $965  
4 Express $102,791           108 $952  
# A tibble: 3 x 4
  Cliente        Segmento Compras Ventas 
  <chr>          <fct>      <int> <chr>  
1 Luis Hernández Maduro        18 $15,744
2 Pedro Cruz     Adulto        13 $14,680
3 Miguel Flores  Adulto        11 $13,742
\end{rsalida}

El Outlet no solo vende más: también tiene el ticket promedio más alto
(\$1,079). Las funciones que preparan la tabla interactiva
(\codigo{tabla\_detalle\_dt()}) las veremos en la
sección~\ref{sec:m6-dt}.

\subsection{La interfaz del dashboard}

Con la lógica probada, pasamos a \archivo{app.R}. Primero paquetes y datos
(fuera del \codigo{server}, una sola vez) y luego encabezado y barra
lateral:

\begin{rnoejecutar}[title=app.R (partes 1 a 3)]
# 1. Paquetes ------------------------------------------------------------
library(shiny)
library(shinydashboard)   # dashboardPage(), box(), valueBox()...
library(dplyr)
library(readr)
library(ggplot2)
library(lubridate)
library(scales)
library(DT)               # tablas interactivas
library(writexl)          # descarga en Excel

# 2. Datos: se leen UNA vez, fuera del server ----------------------------
datos <- cargar_datos_dashboard()
fecha_min <- min(datos$fecha)
fecha_max <- max(datos$fecha)
opciones_tiendas    <- sort(unique(datos$tienda))
opciones_categorias <- sort(unique(datos$categoria))

# 3. Interfaz: encabezado y barra lateral --------------------------------
encabezado <- dashboardHeader(title = "Ventas 2023")

barra_lateral <- dashboardSidebar(
  sidebarMenu(
    id = "menu",
    menuItem("Resumen",   tabName = "resumen",   icon = icon("gauge")),
    menuItem("Productos", tabName = "productos", icon = icon("box")),
    menuItem("Tiendas",   tabName = "tiendas",   icon = icon("store")),
    menuItem("Clientes",  tabName = "clientes",  icon = icon("users")),
    menuItem("Datos",     tabName = "datos",     icon = icon("table"))
  ),
  # Filtros globales: afectan a TODAS las pestañas
  dateRangeInput("fechas", "Periodo:", start = fecha_min, end = fecha_max,
                 min = fecha_min, max = fecha_max, format = "dd/mm/yyyy",
                 language = "es", separator = " a "),
  selectizeInput("tiendas", "Tiendas:", choices = opciones_tiendas,
                 multiple = TRUE, options = list(placeholder = "Todas")),
  selectizeInput("categorias", "Categorías:",
                 choices = opciones_categorias, multiple = TRUE,
                 options = list(placeholder = "Todas")),
  actionButton("limpiar", "Quitar filtros", icon = icon("rotate-left"))
)
\end{rnoejecutar}

Los filtros van \emph{después} del \codigo{sidebarMenu()}, dentro de
\codigo{dashboardSidebar()}: así quedan visibles en todas las pestañas, lo
que deja claro que son globales. Los \codigo{selectizeInput()} múltiples
usan el marcador ``Todas'' para indicar que vacío significa sin filtro.

El cuerpo tiene un \codigo{tabItem()} por cada \codigo{menuItem()}. Cada
fila es un \codigo{fluidRow()} y cada gráfica vive en una \codigo{box()}:

\begin{rnoejecutar}[title=app.R (parte 4): cuerpo]
# 4. Interfaz: cuerpo con una tabItem por pestaña del menú ---------------
cuerpo <- dashboardBody(
  tabItems(
    tabItem(tabName = "resumen",
      fluidRow(                                   # fila de KPIs arriba
        valueBoxOutput("kpi_ventas", width = 3),
        valueBoxOutput("kpi_transacciones", width = 3),
        valueBoxOutput("kpi_ticket", width = 3),
        valueBoxOutput("kpi_clientes", width = 3)
      ),
      fluidRow(
        box(title = "Ventas por mes", width = 8, status = "primary",
            solidHeader = TRUE, plotOutput("graf_evolucion", height = 300)),
        box(title = "Ventas por método de pago", width = 4,
            status = "primary", solidHeader = TRUE,
            plotOutput("graf_pago", height = 300))
      )
    ),
    tabItem(tabName = "productos",
      fluidRow(
        box(title = "Top 10 productos", width = 6, status = "primary",
            plotOutput("graf_top_productos", height = 380)),
        box(title = "Ventas por categoría", width = 6, status = "primary",
            plotOutput("graf_categorias", height = 380))
      )
    ),
    tabItem(tabName = "tiendas",
      fluidRow(
        box(title = "Ventas por tienda", width = 7, status = "primary",
            plotOutput("graf_tiendas", height = 380)),
        box(title = "Indicadores por tienda", width = 5,
            status = "primary", tableOutput("tabla_tiendas"))
      )
    ),
    tabItem(tabName = "clientes",
      fluidRow(
        infoBoxOutput("info_clientes"),
        infoBoxOutput("info_premium"),
        infoBoxOutput("info_frecuencia")
      ),
      fluidRow(
        box(title = "Ventas por segmento", width = 6, status = "primary",
            plotOutput("graf_segmentos", height = 320)),
        box(title = "Top 10 clientes", width = 6, status = "primary",
            tableOutput("tabla_clientes"))
      )
    ),
    tabItem(tabName = "datos",
      fluidRow(
        box(title = "Detalle de transacciones", width = 12,
            status = "primary",
            downloadButton("descargar_csv", "Descargar CSV"),
            downloadButton("descargar_excel", "Descargar Excel"),
            br(), br(),
            DTOutput("tabla_detalle"))
      )
    )
  )
)

ui <- dashboardPage(encabezado, barra_lateral, cuerpo, skin = "blue")
\end{rnoejecutar}

Observa la fila de KPIs: cuatro \codigo{valueBoxOutput()} de ancho 3 suman
exactamente 12 columnas. En la pestaña Clientes, los \codigo{infoBoxOutput()}
tienen el ancho por defecto (4), así que caben tres por fila.

\subsection{El servidor del dashboard}

\begin{rnoejecutar}[title=app.R (parte 5): reactivos, botón y KPIs]
# 5. Servidor: datos filtrados, KPIs y botón para limpiar filtros --------
server <- function(input, output, session) {

  # La pieza central: TODO lo demás depende de datos_filtrados()
  datos_filtrados <- reactive({
    req(input$fechas)
    filtrar_datos(datos, input$fechas, input$tiendas, input$categorias)
  }) %>%
    bindCache(input$fechas, input$tiendas, input$categorias)

  kpis <- reactive(calcular_kpis(datos_filtrados()))

  # Igual que datos_filtrados(), pero si los filtros no dejan ninguna
  # venta, las salidas que lo usan muestran un mensaje amable
  datos_validos <- reactive({
    validate(need(nrow(datos_filtrados()) > 0,
                  "No hay ventas con estos filtros."))
    datos_filtrados()
  })

  observeEvent(input$limpiar, {
    updateDateRangeInput(session, "fechas", start = fecha_min,
                         end = fecha_max)
    updateSelectizeInput(session, "tiendas", selected = character(0))
    updateSelectizeInput(session, "categorias", selected = character(0))
  })

  output$kpi_ventas <- renderValueBox({
    valueBox(formato_pesos(kpis()$ventas), "Ventas totales",
             icon = icon("dollar-sign"), color = "blue")
  })
  output$kpi_transacciones <- renderValueBox({
    valueBox(formato_numero(kpis()$transacciones), "Transacciones",
             icon = icon("receipt"), color = "blue")
  })
  output$kpi_ticket <- renderValueBox({
    valueBox(formato_pesos(kpis()$ticket_promedio), "Ticket promedio",
             icon = icon("cart-shopping"), color = "blue")
  })
  output$kpi_clientes <- renderValueBox({
             icon = icon("users"), color = "blue")
  })
\end{rnoejecutar}

La arquitectura reactiva es la de la figura~\ref{fig:m6-grafo}, un poco más
grande: \codigo{datos\_filtrados()} depende de los tres filtros; de él
cuelgan \codigo{kpis()} (que alimenta los \codigo{valueBox()} y los
\codigo{infoBox()}) y \codigo{datos\_validos()} (que alimenta las gráficas y
muestra un mensaje si no hay ventas). El botón ``Quitar filtros'' es un
\codigo{observeEvent()} que devuelve cada entrada a su valor inicial con las
funciones \codigo{update...()}. Sobre \codigo{bindCache()} hablaremos en la
sección~\ref{sec:m6-buenas-practicas}.

\begin{rnoejecutar}[title=app.R (partes 6 a 8): salidas, descargas y lanzamiento]
  # 6. Servidor: gráficas y tablas de cada pestaña -----------------------
  output$graf_evolucion <- renderPlot({
    grafica_evolucion(datos_validos())
  }, res = 96)
  output$graf_pago <- renderPlot({
    grafica_metodo_pago(datos_validos())
  }, res = 96)
  output$graf_top_productos <- renderPlot({
    grafica_top_productos(datos_validos())
  }, res = 96)
  output$graf_categorias <- renderPlot({
    grafica_categorias(datos_validos())
  }, res = 96)
  output$graf_tiendas <- renderPlot({
    grafica_tiendas(datos_validos())
  }, res = 96)
  output$tabla_tiendas <- renderTable(tabla_tiendas(datos_filtrados()))

  output$info_clientes <- renderInfoBox({
    infoBox("Clientes activos", formato_numero(kpis()$clientes),
            icon = icon("users"), color = "blue")
  })
  output$info_premium <- renderInfoBox({
    infoBox("Clientes premium", formato_pct(kpis()$pct_premium),
            icon = icon("star"), color = "blue")
  })
  output$info_frecuencia <- renderInfoBox({
    infoBox("Compras por cliente",
            formato_numero(kpis()$compras_cliente, decimales = 1),
            icon = icon("repeat"), color = "blue")
  })
  output$graf_segmentos <- renderPlot({
    grafica_segmentos(datos_validos())
  }, res = 96)
  output$tabla_clientes <- renderTable(
    tabla_top_clientes(datos_filtrados()))

  output$tabla_detalle <- renderDT(tabla_detalle_dt(datos_filtrados()))

  # 7. Servidor: descargas (respetan los filtros activos) ----------------
  output$descargar_csv <- downloadHandler(
    filename = function() paste0("ventas_filtradas_", Sys.Date(), ".csv"),
    content  = function(file) {
      write_excel_csv(datos_detalle(datos_filtrados()), file)
    }
  )
  output$descargar_excel <- downloadHandler(
    filename = function() paste0("ventas_filtradas_", Sys.Date(), ".xlsx"),
    content  = function(file) {
      write_xlsx(datos_detalle(datos_filtrados()), file)
    }
  )
}

# 8. Lanzar la app -------------------------------------------------------
shinyApp(ui, server)
\end{rnoejecutar}

Cada \codigo{renderPlot()} tiene una sola línea: llama a su función con los
datos válidos. Si mañana cambias el diseño de una gráfica, lo haces en
\archivo{R/funciones\_dashboard.R} sin tocar \archivo{app.R}.

\subsection{Probar el dashboard completo}

\begin{rcode}
app03 <- shinyAppDir("latex/codigo/apps/03-dashboard-ejecutivo")
testServer(app03, {
  session$setInputs(fechas = as.Date(c("2023-01-01", "2023-12-31")),
                    tiendas = NULL, categorias = NULL)
  cat("Ventas:", kpis()$ventas, "| Clientes:", kpis()$clientes, "\n")

  # El usuario elige dos tiendas y una categoría
  session$setInputs(tiendas = c("Centro", "Norte"),
                    categorias = "Electrónica")
  cat("Ventas:", round(kpis()$ventas, 2),
      "| Transacciones:", kpis()$transacciones, "\n")

  # Filtros sin ventas: las gráficas muestran el mensaje de validate()
  session$setInputs(tiendas = "Express", categorias = "Ropa",
                    fechas = as.Date(c("2023-01-01", "2023-01-31")))
  tryCatch(output$graf_evolucion,
           error = function(e) cat("Gráfica:", conditionMessage(e), "\n"))
  cat("Ticket en el valueBox:", formato_pesos(kpis()$ticket_promedio), "\n")
})
\end{rcode}
\begin{rsalida}
Ventas: 959992.7 | Clientes: 198 
Ventas: 32911.76 | Transacciones: 31 
Gráfica: No hay ventas con estos filtros. 
Ticket en el valueBox: -- 
\end{rsalida}

La prueba confirma los tres escenarios: todo el año (\$959,992.70 y 198
clientes), un filtro combinado (31 transacciones) y un filtro sin ventas,
donde la gráfica muestra el mensaje amable y la tarjeta del ticket muestra
``-{}-'' en lugar de \codigo{NA}.

\begin{hazlotu}
Ejecuta \codigo{shiny::runApp("latex/codigo/apps/03-dashboard-ejecutivo")}.
Recorre las cinco pestañas; luego elige solo la tienda ``Outlet'' y
comprueba que \emph{todas} las pestañas cambian. Pulsa ``Quitar filtros''.
¿Qué método de pago concentra más ventas en el Outlet?
\end{hazlotu}

% ----------------------------------------------------------------------------
\section{Descargas, tablas interactivas y gráficas interactivas}
\label{sec:m6-descargas}

\subsection{Botones de descarga con \texorpdfstring{\codigo{downloadHandler()}}{downloadHandler()}}

Lo primero que pide cualquier usuario de un dashboard es ``¿me lo puedes
pasar en Excel?''. Un botón de descarga tiene, como toda salida, dos mitades:
\codigo{downloadButton("id", "Texto")} en la \codigo{ui} y
\codigo{output\$id <- downloadHandler(filename, content)} en el
\codigo{server}:

\begin{itemize}
  \item \codigo{filename}: una \emph{función} que devuelve el nombre del
        archivo. Incluir la fecha (\codigo{Sys.Date()}) evita que el usuario
        sobrescriba descargas anteriores.
  \item \codigo{content}: una \emph{función} que recibe una ruta temporal
        (\codigo{file}) y escribe ahí el archivo. Shiny se encarga de
        enviarlo al navegador.
\end{itemize}

En la app 03 (parte 7 del servidor) hay dos botones, uno para CSV y otro
para Excel, y ambos descargan \codigo{datos\_detalle(datos\_filtrados())}:
\textbf{respetan los filtros activos}, que es lo que el usuario espera.
Para el CSV usamos \codigo{write\_excel\_csv()} de \paquete{readr} en lugar
de \codigo{write\_csv()}: agrega una marca invisible al inicio del archivo
para que Excel muestre bien los acentos (``Categoría'', ``Tarjeta
Crédito''). Para Excel usamos \codigo{write\_xlsx()} de \paquete{writexl}.

Como \codigo{content} es una función normal, puedes probar lo que hará
escribiendo tú mismo los archivos:

\begin{rcode}
library(writexl)
dir.create("resultados", showWarnings = FALSE)
detalle <- datos_detalle(filtrar_datos(datos, anio, tiendas = "Outlet"))
write_excel_csv(detalle, "resultados/ventas_outlet.csv")
write_xlsx(detalle, "resultados/ventas_outlet.xlsx")
dim(detalle)
names(detalle)
\end{rcode}
\begin{rsalida}
[1] 106   9
[1] "Fecha"          "Tienda"         "Categoría"     
[4] "Producto"       "Cliente"        "Cantidad"      
[7] "Precio"         "Total"          "Método de pago"
\end{rsalida}

Las 106 transacciones del Outlet, con 9 columnas de nombres legibles, listas
para quien las pida. En la prueba con \codigo{testServer()} también puedes
leer \codigo{output\$descargar\_csv}: devuelve la ruta del archivo generado,
que puedes abrir con \codigo{read\_csv()} para comprobar su contenido.

\subsection{Tablas interactivas con \texorpdfstring{\paquete{DT}}{DT}}
\label{sec:m6-dt}

\codigo{renderTable()} está bien para tablas pequeñas (un top 10). Para
mostrar cientos o miles de filas necesitas búsqueda, orden por columna y
paginación: eso hace \paquete{DT}, un puente a la biblioteca DataTables de
JavaScript. La pareja es \codigo{DTOutput()} / \codigo{renderDT()}, y
dentro se usa \codigo{datatable()}:

\begin{rcode}
# Textos de la tabla DT en español (sin depender de internet)
idioma_dt <- list(
  search = "Buscar:", lengthMenu = "Mostrar _MENU_ filas",
  info = "Filas _START_ a _END_ de _TOTAL_",
  infoEmpty = "Sin resultados", zeroRecords = "No hay coincidencias",
  paginate = list(previous = "Anterior", `next` = "Siguiente")
)

tabla_detalle_dt <- function(datos) {
  datatable(datos_detalle(datos), rownames = FALSE,
            options = list(pageLength = 10, language = idioma_dt)) %>%
    formatCurrency(c("Precio", "Total"), currency = "$", digits = 2)
}

tabla_dt <- tabla_detalle_dt(datos)
class(tabla_dt)
\end{rcode}
\begin{rsalida}
[1] "datatables" "htmlwidget"
\end{rsalida}

El resultado es un \termino{htmlwidget}: un objeto de R que se dibuja con
JavaScript en el navegador (en RStudio, escribe \codigo{tabla\_dt} en la
consola y lo verás en el visor). Las opciones más útiles:

\begin{itemize}
  \item \codigo{formatCurrency(columnas, currency = "\$")}: formato de moneda
        sin convertir los números a texto, así que ordenar por esa columna
        sigue funcionando bien. También existen \codigo{formatPercentage()},
        \codigo{formatRound()} y \codigo{formatDate()}.
  \item \codigo{options = list(pageLength = 10)}: filas por página.
  \item \codigo{language = idioma\_dt}: traduce ``Search'', ``Next'', etc. Lo
        definimos a mano para no depender de internet.
  \item \codigo{filter = "top"}: agrega una caja de búsqueda encima de cada
        columna.
  \item \codigo{rownames = FALSE}: oculta la columna de números de fila.
\end{itemize}

\begin{advertencia}
\paquete{shiny} también tiene funciones llamadas \codigo{dataTableOutput()}
y \codigo{renderDataTable()}, antiguas y distintas a las de \paquete{DT}. Por
eso al cargar \paquete{DT} aparece el aviso de que las ``enmascara''. Usa
siempre \codigo{DT::DTOutput()} y \codigo{DT::renderDT()} y evitarás
confusiones.
\end{advertencia}

\subsection{Gráficas interactivas con \texorpdfstring{\paquete{plotly}}{plotly}}

Con \paquete{plotly} una gráfica de \paquete{ggplot2} gana tooltips (al
pasar el ratón ves el valor exacto), zoom y la opción de ocultar series. La
forma más rápida es \codigo{ggplotly()}, que convierte un ggplot ya hecho:

\begin{rcode}
library(plotly)
graf_interactiva <- ggplotly(p_tiendas)
class(graf_interactiva)
\end{rcode}
\begin{rsalida}

Attaching package: ‘plotly’

The following object is masked from ‘package:ggplot2’:

    last_plot

The following object is masked from ‘package:stats’:

    filter

The following object is masked from ‘package:graphics’:

    layout

[1] "plotly"     "htmlwidget"
\end{rsalida}

En la app basta con cambiar la pareja de funciones:

\begin{rnoejecutar}
# En la ui, en lugar de plotOutput("graf_tiendas"):
plotly::plotlyOutput("graf_tiendas", height = 380)

# En el server, en lugar de renderPlot():
output$graf_tiendas <- plotly::renderPlotly({
  ggplotly(grafica_tiendas(datos_validos()))
})
\end{rnoejecutar}

\begin{consejo}
No conviertas todo a plotly ``porque se ve moderno''. La interactividad
ayuda en gráficas con muchos puntos o cuando el usuario necesita el valor
exacto; en un KPI o un ranking de 10 barras solo añade ruido y tiempo de
carga. Las etiquetas de \codigo{geom\_text()} y algunos detalles de estilo
no siempre se traducen igual; revisa el resultado.
\end{consejo}

% ----------------------------------------------------------------------------
\section{Buenas prácticas profesionales}
\label{sec:m6-buenas-practicas}

\subsection{Separa la lógica de la interfaz}

Ya lo aplicaste en las dos apps: las funciones de datos, KPIs y gráficas no
tocan \codigo{input} ni \codigo{output}. Regla práctica: \emph{si un bloque
dentro de un \codigo{render...()} tiene más de 3 o 4 líneas, conviértelo en
una función} y guárdala en \archivo{R/}. Ganas pruebas sencillas,
reutilización (las mismas funciones sirven para el reporte automático del
Módulo~\ref{cap:reportes}) y un servidor que se lee como un índice.

\subsection{Lee los datos una sola vez y precalcula}

Todo lo que está fuera del \codigo{server} se ejecuta al arrancar y lo
comparten todos los usuarios; lo que está dentro se ejecuta \emph{por cada
usuario}. Por eso:

\begin{itemize}
  \item Lee y une los datos fuera del \codigo{server} (como en ambas apps).
  \item Si los datos son grandes, prepáralos antes con un script aparte y
        guárdalos ya resumidos (por ejemplo en \archivo{.rds} con
        \codigo{saveRDS()}, que se lee mucho más rápido que un CSV), o
        consúltalos en una base de datos agregando en SQL
        (Módulo~\ref{cap:bases-datos}).
  \item Usa \codigo{reactive()} para no repetir cálculos entre salidas.
\end{itemize}

\subsection{Caché con \texorpdfstring{\codigo{bindCache()}}{bindCache()}}

\codigo{reactive()} recuerda \emph{el último} resultado. \codigo{bindCache()}
va más allá: recuerda los resultados para \emph{cada combinación} de valores
que le indiques, y además los comparte entre usuarios. En la app 03:

\begin{rnoejecutar}
datos_filtrados <- reactive({
  req(input$fechas)
  filtrar_datos(datos, input$fechas, input$tiendas, input$categorias)
}) %>%
  bindCache(input$fechas, input$tiendas, input$categorias)
\end{rnoejecutar}

Si la directora mira ``Outlet'', luego ``Centro'' y vuelve a ``Outlet'', la
tercera vez el resultado sale de la caché sin recalcular; y si un gerente
abre la app con los mismos filtros, también lo aprovecha. También funciona
con \codigo{renderPlot()}: \codigo{renderPlot(...) \%>\%
bindCache(input\$tiendas)} guarda las imágenes ya dibujadas. La regla de oro:
en \codigo{bindCache()} pon \emph{todas} las entradas de las que depende el
cálculo; si olvidas una, mostrarás resultados viejos.

\subsection{Módulos de Shiny: componentes reutilizables}

Cuando un mismo bloque (una tarjeta de KPI, un selector de periodo, una
gráfica con su tabla) se repite varias veces, copiar y pegar lleva a
errores: hay que inventar ids distintos (\codigo{kpi1}, \codigo{kpi2}\ldots)
y cambiar cada copia a mano. Un \termino{módulo de Shiny} es un par de
funciones (una de \codigo{ui} y otra de \codigo{server}) que encapsulan ese
bloque. El truco está en el \termino{espacio de nombres}
(\emph{namespace}): \codigo{NS(id)} antepone el id del módulo a cada id
interno, de modo que dos copias nunca chocan.

\begin{ejemplo}{Módulo ``tarjeta KPI''}
Queremos una tarjeta reutilizable que reciba un valor reactivo, un título,
un ícono y una función de formato, y muestre el \codigo{valueBox()}.
\end{ejemplo}

\begin{rcode}
# Parte visual del módulo: todos los ids pasan por ns()
tarjeta_kpi_ui <- function(id, ancho = 3) {
  ns <- NS(id)
  valueBoxOutput(ns("caja"), width = ancho)
}

# Parte lógica: recibe un reactivo 'valor' desde fuera
tarjeta_kpi_server <- function(id, valor, titulo, icono,
                               formato = formato_pesos) {
  moduleServer(id, function(input, output, session) {
    output$caja <- renderValueBox({
      valueBox(formato(valor()), titulo, icon = icon(icono),
               color = "blue")
    })
  })
}

tarjeta_kpi_ui("ventas")          # el id interno se vuelve "ventas-caja"
\end{rcode}
\begin{rsalida}
<div class="shiny-html-output col-sm-3" id="ventas-caja"></div>
\end{rsalida}

El \codigo{id="ventas-caja"} muestra el namespace en acción. Con el módulo,
la fila de KPIs del dashboard quedaría así:

\begin{rnoejecutar}
# ui
fluidRow(
  tarjeta_kpi_ui("ventas"),
  tarjeta_kpi_ui("ticket"),
  tarjeta_kpi_ui("clientes")
)
# server
tarjeta_kpi_server("ventas", reactive(kpis()$ventas),
                   "Ventas totales", "dollar-sign")
tarjeta_kpi_server("ticket", reactive(kpis()$ticket_promedio),
                   "Ticket promedio", "cart-shopping")
tarjeta_kpi_server("clientes", reactive(kpis()$clientes),
                   "Clientes activos", "users", formato = formato_numero)
\end{rnoejecutar}

Los módulos también se prueban con \codigo{testServer()}; los argumentos
del módulo se pasan con \codigo{args}:

\begin{rcode}
testServer(tarjeta_kpi_server,
           args = list(valor = reactive(kpis_total$clientes),
                       titulo = "Clientes activos", icono = "users",
                       formato = formato_numero), {
  cat(as.character(output$caja$html))
})
\end{rcode}
\begin{rsalida}
<div class="small-box bg-blue">
  <div class="inner">
    <h3>198</h3>
    <p>Clientes activos</p>
  </div>
  <div class="icon-large">
    <i class="fas fa-users" role="presentation" aria-label="users icon"></i>
  </div>
</div>
\end{rsalida}

\subsection{Prueba tu app con \texorpdfstring{\codigo{testServer()}}{testServer()}}

Ya lo usaste en todo el módulo. En un proyecto real, guarda esas pruebas en
una carpeta \archivo{tests/} y ejecútalas antes de publicar cada cambio (el
paquete \paquete{testthat} las organiza con \codigo{expect\_equal()}). Para
probar además la parte visual (clics reales en un navegador) existe el
paquete \paquete{shinytest2}.

\subsection{Estructura de carpetas de un proyecto Shiny}

Esta es una organización que escala bien de una app pequeña a una grande:

\begin{textocodigo}
dashboard-ventas/
|-- app.R                    <- solo ui, server y shinyApp()
|-- R/                       <- Shiny carga estos archivos solo
|   |-- funciones_datos.R    <- cargar, limpiar, filtrar
|   |-- funciones_graficas.R <- una función por gráfica
|   +-- mod_tarjeta_kpi.R    <- módulos (ui + server)
|-- datos/                   <- datos ya preparados (.rds, .csv)
|-- www/                     <- logo, CSS, imágenes
|-- tests/
|   +-- testthat/
|       +-- test-kpis.R      <- pruebas con testServer()
+-- preparar_datos.R         <- script que genera datos/ (no se publica)
\end{textocodigo}

La carpeta \archivo{www/} es especial: lo que pongas ahí (por ejemplo
\archivo{logo.png}) queda disponible en la página como
\codigo{img(src = "logo.png")}.

% ----------------------------------------------------------------------------
\section{Publicar tu dashboard}
\label{sec:m6-publicar}

Mientras la app corre en tu computadora, solo tú la ves. Para que la use tu
equipo, debe correr en un \termino{servidor}. Opciones:

\begin{itemize}
  \item \textbf{shinyapps.io}: el servicio en la nube de Posit (la empresa
        detrás de RStudio). Tiene un plan gratuito con horas limitadas al
        mes, ideal para prototipos y portafolio.
  \item \textbf{Posit Connect}: servidor empresarial de pago que se instala
        en la infraestructura de la empresa; incluye autenticación con las
        cuentas corporativas, control de acceso por app y publicación de
        reportes R Markdown y Quarto.
  \item \textbf{Shiny Server} (versión gratuita, de código abierto): lo
        instala el área de TI en un servidor Linux propio. No incluye
        autenticación; hay que ponerla con otras herramientas.
\end{itemize}

Para publicar en shinyapps.io, crea una cuenta en
\url{https://www.shinyapps.io}, copia tu \emph{token} (menú
\menu{Account > Tokens}) y ejecuta una sola vez:

\begin{rnoejecutar}
install.packages("rsconnect")
rsconnect::setAccountInfo(name   = "tu_usuario",
                          token  = "TOKEN_DE_EJEMPLO",
                          secret = "SECRETO_DE_EJEMPLO")
\end{rnoejecutar}

Antes de publicar, \textbf{copia los datos dentro de la carpeta de la app}:
en el servidor no existe la carpeta del curso, así que
\codigo{buscar\_datasets()} no encontraría nada. Basta con crear
\archivo{datasets/} dentro de la carpeta de la app (la función busca
primero ahí, en el nivel 0). Luego:

\begin{rnoejecutar}
rsconnect::deployApp("latex/codigo/apps/03-dashboard-ejecutivo",
                     appName = "dashboard-ventas")
\end{rnoejecutar}

En uno o dos minutos obtienes una dirección del tipo
\url{https://tu_usuario.shinyapps.io/dashboard-ventas/}. Cada vez que
cambies la app, vuelve a ejecutar \codigo{deployApp()}. Si falla al
arrancar en el servidor, revisa el registro con
\codigo{rsconnect::showLogs()}: casi siempre es un paquete no instalado o
un archivo de datos que no se subió.

\begin{advertencia}
\textbf{Seguridad.} Una app en shinyapps.io con plan gratuito es
\emph{pública}: cualquiera con la dirección la ve. Nunca publiques así datos
de clientes (nombres, correos, teléfonos), salarios o cifras
confidenciales. Para datos internos usa un servidor con autenticación
(Posit Connect, planes de pago de shinyapps.io o un servidor de la empresa),
publica solo datos agregados o anonimizados, y jamás escribas contraseñas
de bases de datos en el código: guárdalas en variables de entorno
(\archivo{.Renviron}), como viste en el Módulo~\ref{cap:bases-datos}.
\end{advertencia}

\subsection{Alternativas sin servidor: flexdashboard y Quarto dashboards}

Si tu tablero no necesita recalcular nada al vuelo (por ejemplo, un reporte
semanal fijo), no necesitas Shiny ni un servidor: \paquete{flexdashboard}
(con R Markdown) y los \emph{dashboards de Quarto} generan un archivo HTML
con aspecto de tablero, gráficas \paquete{plotly} y tablas \paquete{DT}
interactivas, que puedes enviar por correo o subir a una intranet. El
encabezado de un dashboard de Quarto se ve así:

\begin{textocodigo}
---
title: "Ventas 2023"
format: dashboard
---
\end{textocodigo}

Verás cómo automatizar este tipo de reportes en el
Módulo~\ref{cap:reportes}. Regla rápida: si el usuario necesita
\emph{filtrar y recalcular}, Shiny; si solo necesita \emph{consultar} un
resultado fijo con algo de interactividad, un HTML estático basta.

% ----------------------------------------------------------------------------
\section{Errores comunes}
\label{sec:m6-errores}

\textbf{1. Usar un reactivo sin paréntesis.} Es el error más famoso. En
esta mini app, \codigo{datos} es un reactivo y lo usamos como si fuera un
data frame:

\begin{rcode}
server_error <- function(input, output, session) {
  datos <- reactive(filter(transacciones, tienda_id == input$tienda))
  output$total <- renderText(sum(datos$total_transaccion))  # ¡error!
}
testServer(server_error, {
  session$setInputs(tienda = 1)
  tryCatch(output$total,
           error = function(e) cat("Error:", conditionMessage(e), "\n"))
})
\end{rcode}
\begin{rsalida}
Error: objeto de tipo 'closure' no es subconjunto 
\end{rsalida}

Con R en inglés el mensaje es \codigo{object of type 'closure' is not
subsettable}. La corrección: \codigo{sum(datos()\$total\_transaccion)}.
Ojo: \codigo{nrow(datos)} sin paréntesis \emph{no} da error, devuelve
\codigo{NULL} y la salida aparece vacía, lo que es aún más desconcertante.

\textbf{2. Olvidar una coma entre elementos de la \codigo{ui}.} R no puede
ni leer el archivo:

\begin{rcode}
codigo_ui <- 'fluidPage(
  titlePanel("Ventas")
  sliderInput("n", "Top N:", 5, 20, 10)
)'
tryCatch(parse(text = codigo_ui),
         error = function(e) cat(conditionMessage(e), "\n"))
\end{rcode}
\begin{rsalida}
<text>:3:3: unexpected symbol
2:   titlePanel("Ventas")
3:   sliderInput
     ^ 
\end{rsalida}

El error señala la línea 3, pero la coma falta al final de la línea 2. Con
una coma \emph{de más} después del último elemento el mensaje suele ser
\codigo{argument is missing, with no default}.

\begin{table}[H]
\centering
\caption{Otros errores frecuentes en Shiny y cómo resolverlos.}
\label{tab:m6-errores}
\small
\begin{tabularx}{\textwidth}{>{\raggedright\arraybackslash}p{4.8cm}
                             >{\raggedright\arraybackslash}X}
\toprule
\textbf{Síntoma o mensaje} & \textbf{Causa y solución} \\
\midrule
\codigo{Operation not allowed without an active reactive context} &
Llamaste a un reactivo (\codigo{datos()}) en la consola o en el cuerpo del
\codigo{server}. Muévelo dentro de un \codigo{render...()},
\codigo{reactive()} u \codigo{observe()}. \\
\codigo{Can't access reactive value 'x' outside of reactive consumer} &
Usaste \codigo{input\$x} fuera de un contexto reactivo. Envuélvelo en
\codigo{reactive()}. \\
\codigo{'datasets/...csv' does not exist in current working directory} &
La app corre con su carpeta como directorio de trabajo. Usa
\codigo{buscar\_datasets()} o copia los datos a la carpeta de la app. \\
Una salida no aparece, sin error & El id de \codigo{xxxOutput("id")} no
coincide con \codigo{output\$id}, o la pareja no corresponde
(\codigo{plotOutput} con \codigo{renderTable}). Revisa la
tabla~\ref{tab:m6-salidas}. \\
Dos controles o salidas dejan de funcionar (en la consola del navegador,
\tecla{F12}, aparece \codigo{Duplicate binding for ID ...}) & Usaste el
mismo id dos veces. Los ids deben ser únicos en toda la app; con módulos,
usa \codigo{ns()}. \\
La gráfica muestra un error rojo al abrir la app & Una entrada aún no tiene
valor. Agrega \codigo{req(input\$...)} al inicio del \codigo{render}. \\
Un filtro múltiple vacío deja la app sin datos & \codigo{input\$x} vale
\codigo{NULL}; revisa \codigo{length(input\$x) > 0} antes de filtrar. \\
\codigo{there is no package called 'DT'} al publicar & El servidor no
tiene el paquete: cárgalo con \codigo{library()} en \archivo{app.R} (no con
\codigo{require()} dentro de funciones) para que \codigo{rsconnect} lo
detecte. \\
La app funciona en tu computadora pero no en el servidor & Rutas absolutas
(\codigo{C:/Users/...}), datos no subidos o paquetes faltantes. Revisa
\codigo{rsconnect::showLogs()}. \\
\bottomrule
\end{tabularx}
\end{table}

\begin{consejo}
Para depurar, pon \codigo{browser()} dentro de un \codigo{render...()} o
\codigo{reactive()}: al ejecutarse, la app se pausa y la consola de RStudio
te deja inspeccionar \codigo{input\$...} y los objetos en ese momento. Escribe
\codigo{c} para continuar. Quítalo antes de publicar.
\end{consejo}

% ----------------------------------------------------------------------------
\begin{resumen}
\begin{itemize}
  \item Un buen dashboard tiene una audiencia y unas preguntas definidas,
        KPIs arriba, filtros juntos y pocas gráficas bien diseñadas. Shiny
        conviene cuando necesitas lógica a la medida en R; Power BI o
        Tableau, para reportes estándar de autoservicio.
  \item Toda app tiene \codigo{ui} (lo que se ve), \codigo{server} (la
        lógica) y \codigo{shinyApp()}. Cada salida es una pareja
        \codigo{xxxOutput()} + \codigo{renderXxx()} con el mismo id.
  \item La reactividad recalcula solo lo que depende de lo que cambió.
        \codigo{reactive()} comparte cálculos; \codigo{observeEvent()} hace
        acciones; \codigo{eventReactive()} espera a un botón; \codigo{req()}
        y \codigo{validate()} evitan errores; \codigo{isolate()} lee sin
        depender. Los reactivos se llaman con paréntesis.
  \item Organiza la pantalla con la cuadrícula de 12 columnas,
        \codigo{sidebarLayout()}, pestañas o \paquete{shinydashboard}
        (\codigo{box()}, \codigo{valueBox()}, \codigo{menuItem()}).
  \item Separa la lógica en funciones (carpeta \archivo{R/}), lee los datos
        una vez, usa \codigo{bindCache()} y módulos, y prueba con
        \codigo{testServer()}.
  \item Publica en shinyapps.io, Posit Connect o Shiny Server, sin exponer
        datos sensibles.
\end{itemize}

\medskip
\small
\begin{tabularx}{\linewidth}{>{\raggedright\arraybackslash}p{4.6cm}X}
\toprule
\textbf{Función} & \textbf{Para qué sirve} \\
\midrule
\codigo{fluidPage()}, \codigo{sidebarLayout()} & Página y layout con barra lateral \\
\codigo{fluidRow()}, \codigo{column()} & Cuadrícula de 12 columnas \\
\codigo{tabsetPanel()}, \codigo{navbarPage()} & Pestañas y barra de navegación \\
\codigo{selectInput()}, \codigo{selectizeInput()} & Menús de una o varias opciones \\
\codigo{dateRangeInput()}, \codigo{sliderInput()} & Periodos y rangos \\
\codigo{actionButton()}, \codigo{downloadButton()} & Botones de acción y de descarga \\
\codigo{renderPlot()}, \codigo{renderText()}, \codigo{renderTable()} & Salidas básicas \\
\codigo{DT::renderDT()}, \codigo{plotly::renderPlotly()} & Tabla y gráfica interactivas \\
\codigo{reactive()}, \codigo{eventReactive()} & Cálculos reactivos compartidos \\
\codigo{observeEvent()}, \codigo{update...Input()} & Acciones y cambios de entradas \\
\codigo{req()}, \codigo{validate(need())} & Condiciones antes de calcular \\
\codigo{dashboardPage()}, \codigo{box()}, \codigo{valueBox()} & Estructura de shinydashboard \\
\codigo{downloadHandler()} & Contenido de una descarga \\
\codigo{bindCache()} & Caché de resultados \\
\codigo{NS()}, \codigo{moduleServer()} & Módulos reutilizables \\
\codigo{testServer()} & Probar el servidor sin navegador \\
\codigo{runApp()}, \codigo{rsconnect::deployApp()} & Ejecutar y publicar \\
\bottomrule
\end{tabularx}
\end{resumen}

% ----------------------------------------------------------------------------
\section{Ejercicios}
\label{sec:m6-ejercicios}

Para los ejercicios que modifican una app, copia primero su carpeta (por
ejemplo, \archivo{latex/codigo/apps/02-explorador-ventas} a
\archivo{mis-apps/02-mi-explorador}) y trabaja sobre la copia. Además de
abrirla en el navegador, comprueba la lógica con \codigo{testServer()}.

\begin{ejercicio}{m6-1}{Calculadora con cantidad e IVA}
\nivel{Básico}
Amplía la app 01: agrega un \codigo{numericInput()} para la
\emph{cantidad} de piezas (valor inicial 1) y un \codigo{checkboxInput()}
``Incluir IVA (16\%)''. La app debe mostrar el total a pagar
(precio $\times$ cantidad, menos el descuento y, si se marca la casilla,
más IVA). Comprueba con \codigo{testServer()} que 3 piezas de \$500 con
10\% de descuento e IVA dan \$1,566.
\end{ejercicio}

\begin{ejercicio}{m6-2}{Filtro por método de pago}
\nivel{Básico}
En el explorador de ventas (app 02) agrega un
\codigo{checkboxGroupInput()} con los cuatro métodos de pago, todos
marcados al inicio, y haz que filtre todas las salidas. ¿Cuántas
transacciones y cuántas ventas quedan si solo dejas ``Efectivo''?
\end{ejercicio}

\begin{ejercicio}{m6-3}{Un KPI nuevo: margen bruto}
\nivel{Intermedio}
En el dashboard ejecutivo (app 03) agrega en la pestaña Resumen una quinta
fila con un \codigo{valueBox()} de \emph{margen bruto}
(ventas $-$ costo $\times$ cantidad) y otro con el \emph{\% de margen}. La
caja del margen debe ser verde si es positivo y roja si es negativo.
Escribe primero la función \codigo{calcular\_margen(datos)} y pruébala: ¿qué
margen tiene la categoría Belleza? (Recuerda el hallazgo de calidad de datos
del curso: hay productos cuyo costo supera su precio).
\end{ejercicio}

\begin{ejercicio}{m6-4}{Aplicar filtros solo al pulsar un botón}
\nivel{Intermedio}
Escribe una app con un \codigo{selectizeInput()} de tiendas (múltiple), un
\codigo{dateRangeInput()} y un botón ``Aplicar''. Las ventas totales y el
número de transacciones deben actualizarse \emph{solo} al pulsar el botón
(usa \codigo{eventReactive()}). Demuestra con \codigo{testServer()} que
cambiar los filtros sin pulsar no altera el resultado.
\end{ejercicio}

\begin{ejercicio}{m6-5}{Una pestaña nueva: días de la semana}
\nivel{Intermedio}
Agrega al dashboard ejecutivo una pestaña ``Días'' (con su
\codigo{menuItem()} y su \codigo{tabItem()}) con una gráfica de barras de
ventas por día de la semana (lunes a domingo, en ese orden) que resalte el
mejor día. Respeta los filtros globales. ¿Qué día vende más en todo 2023?
\end{ejercicio}

\begin{ejercicio}{m6-6}{Descargar el top de productos}
\nivel{Intermedio}
En la pestaña ``Top productos'' de la app 02 agrega dos botones: ``Descargar
CSV'' y ``Descargar Excel''. El archivo debe contener exactamente la tabla
que se está viendo (con el número de productos del deslizador y los filtros
activos) y llamarse \archivo{top\_productos\_AAAA-MM-DD}. Prueba el
\codigo{downloadHandler()} con \codigo{testServer()}.
\end{ejercicio}

\begin{ejercicio}{m6-7}{Un módulo de ranking reutilizable}
\nivel{Avanzado}
Crea un módulo \codigo{ranking\_ui(id, titulo)} /
\codigo{ranking\_server(id, datos, columna)} que muestre dentro de una
\codigo{box()} un \codigo{sliderInput()} ``Mostrar top'' (3 a 10) y una
gráfica de barras horizontales con las ventas por la columna indicada.
Úsalo dos veces en la misma app: para productos
(\codigo{"nombre\_producto"}) y para ciudades (\codigo{"ciudad"}), y verifica
con \codigo{testServer()} que cada copia responde a su propio deslizador.
\end{ejercicio}

\begin{ejercicio}{m6-8}{Reto: dashboard de Recursos Humanos}
\nivel{Reto}
Construye desde cero, con \paquete{shinydashboard}, un dashboard para el
área de Recursos Humanos con \archivo{empleados.csv}: filtros de
departamento y sucursal (múltiples); cuatro \codigo{valueBox()} (plantilla,
salario mensual promedio, antigüedad promedio y \% con maestría o
doctorado); una gráfica de salario promedio por puesto; una tabla
\paquete{DT} con formato de moneda y un botón para descargar la lista en
Excel. Separa la lógica en funciones y pruébalas.
\end{ejercicio}
